% =============================================================================
%  Texto, tempo e bits
%  Cobre os subprojetos strings, chrono, bits e file_io.
% =============================================================================

% =============================================================================
\section{Texto, tempo e bits}
\label{sec:texto-tempo-bits}
% =============================================================================

Três assuntos que parecem não ter nada em comum ocupam este capítulo: cadeias de
caracteres, medição de tempo e manipulação de bits. Eles estão juntos por uma
razão precisa --- são os três lugares em que o programador acredita saber o que
está acontecendo e \textbf{está errado}. Concatenar duas \emph{strings} parece
grátis e aloca memória quatro vezes; cronometrar um trecho parece trivial e
mede o relógio em vez do código; inverter os bits de um byte parece produzir
outro byte e produz um inteiro negativo de 32 bits.

O capítulo percorre os subprojetos \projeto{strings}, \projeto{chrono},
\projeto{bits} e \projeto{file\_io}. O último é a saída natural dos três
primeiros: texto formatado, instantes registrados e bytes serializados
precisam, no fim, chegar a um arquivo.

\begin{porque}[por que estes três assuntos merecem medição, e não afirmação]
Em quase todo o resto do C++ o erro é detectável: o compilador recusa, o
\emph{sanitizer} aborta, o teste falha. Nos três assuntos deste capítulo o erro
é \textbf{quantitativo}. Um programa que aloca quatro vezes onde poderia alocar
uma vez continua correto; um \emph{benchmark} que mede o compilador em vez do
algoritmo continua produzindo números; uma \texttt{regex} que custa cem vezes
mais que um \texttt{find} continua devolvendo a resposta certa. A única forma
de descobrir a diferença é instrumentar e medir --- e é isso que os subprojetos
fazem, contando alocações de verdade e cronometrando com o relógio certo.
\end{porque}

Vale fixar de saída um vocabulário mínimo, porque ele reaparece em todas as
seções.

\begin{description}
\item[Alocação.] O pedido de memória ao sistema de tempo de execução, feito por
\texttt{operator new} ou \texttt{malloc}. Ela custa muito mais que uma
atribuição: envolve procurar um bloco livre, atualizar metadados e,
eventualmente, pedir mais memória ao sistema operacional. Contar alocações é,
por isso, uma métrica melhor que cronometrar --- é determinística.
\item[\emph{Heap}.] A região de memória de onde vêm as alocações dinâmicas, em
oposição à \emph{pilha}, que é a região onde vivem as variáveis locais e cujo
espaço é reservado e devolvido pelo simples avanço e recuo de um registrador.
\item[Linha de cache.] A unidade mínima que o processador traz da memória
principal para a memória rápida interna --- tipicamente 64 bytes. Ler um byte
custa o mesmo que ler os 64 da linha; por isso dados vizinhos são quase
gratuitos e dados espalhados custam uma espera por elemento.
\item[Unidade de tradução.] Um arquivo \texttt{.cpp} depois de o
pré-processador ter inserido todos os cabeçalhos que ele inclui. É o que o
compilador vê de cada vez, e a razão de o subprojeto \projeto{strings} precisar
declarar o \texttt{operator new} instrumentado em uma unidade só.
\end{description}

% =============================================================================
\subsection{Texto}
\label{sec:texto}
% =============================================================================

O subprojeto \projeto{strings} abre com três números que resumem o problema
inteiro:

\begin{lstlisting}[style=bt, language=bash,
    caption={Saída de \texttt{src/strings/}: o resumo do subprojeto.}]
strings de até 15 chars          →  0 alocações (SSO)
200 chars, um a um, sem reserve  →  4 alocações
extrair um campo de log:
  find + from_chars    →   0 ms
  std::regex           → 133 ms      ← mesmo com a regex fora do laço
\end{lstlisting}

As três linhas dizem respectivamente que texto curto é barato, que texto montado
sem planejamento é caro, e que a ferramenta genérica de casamento de padrões
custa duas ordens de grandeza mais que o código específico. As subseções
seguintes explicam cada uma.

\subsubsection{\classe{std::string} por dentro: a SSO, medida}

O primeiro problema que \classe{std::string} resolve é a posse: quem libera a
memória de uma sequência de caracteres de tamanho variável. A resposta é
\gls{raii} --- o construtor aloca, o destrutor libera --- e ela custa uma
alocação por \emph{string}. Para nomes, chaves de dicionário e identificadores,
que são curtos e existem aos milhares, essa alocação seria ruinosa.

A solução das implementações é a \textbf{SSO}, sigla de \emph{Small String
Optimization}: o objeto \texttt{string} reserva alguns bytes dentro de si mesmo
e guarda ali as \emph{strings} curtas, sem tocar na \emph{heap}. O detalhe é que
o padrão não obriga nem documenta o limite --- ele é uma decisão de
implementação. O subprojeto, então, mede:

\begin{lstlisting}[style=bt, language=bash,
    caption={Saída de \texttt{src/strings/}: o limite da SSO na libstdc++.}]
sizeof(std::string) = 32 bytes
 1 chars → 0 alocações  ← cabe no SSO
15 chars → 0 alocações  ← cabe no SSO
16 chars → 1 alocação
\end{lstlisting}

A medição é feita substituindo o \texttt{operator new} global por uma versão que
incrementa um contador --- a mesma técnica de diagnóstico usada em
\cref{sec:vocabulary}. E há uma segunda prova, estrutural: numa \emph{string}
curta, o endereço devolvido por \texttt{.data()} fica 16 bytes depois do
endereço do próprio objeto, ou seja, \emph{dentro} dele; numa \emph{string}
longa, fica a uma distância arbitrária, na \emph{heap}.

Três consequências práticas decorrem disso, e a segunda contraria a intuição
corrente:

\begin{itemize}
\item \emph{strings} curtas são quase de graça --- não evite \classe{std::string}
      em nomes, chaves e identificadores por medo de alocação;
\item mas o \gls{move} de uma \emph{string} curta \textbf{copia} os bytes: não há
      ponteiro a roubar, porque os dados estão dentro do objeto. O \emph{move} de
      \texttt{string} só é $O(1)$ \emph{acima} do limite da SSO;
\item \texttt{sizeof(std::string)} é 32 bytes, então um
      \texttt{vector<string>} de um milhão de elementos gasta 32~MB apenas de
      cabeçalhos, antes de qualquer caractere.
\end{itemize}

\begin{armadilha}[o \emph{move} barato de \texttt{string} não é sempre barato]
A regra de bolso ``mover é $O(1)$, copiar é $O(n)$'' vale para
\classe{std::vector} sem exceção e para \classe{std::string} apenas acima de 15
caracteres. Abaixo disso, mover e copiar custam exatamente o mesmo --- uma cópia
de 16 bytes. Isso não é motivo para deixar de mover; é motivo para não
\emph{esperar} ganho ao mover uma coleção de \emph{strings} curtas, e para
desconfiar de um perfil que mostra \texttt{std::string} caro sem que haja
alocação envolvida.
\end{armadilha}

\subsubsection{\texttt{capacity}, crescimento e \texttt{reserve}}

\texttt{size()} é quantos caracteres existem; \texttt{capacity()} é quantos
cabem antes de a \emph{string} precisar de um \emph{buffer} maior. Quando a
capacidade acaba, a implementação aloca um bloco maior, copia o conteúdo e
libera o antigo --- e a política de crescimento é a mesma do \texttt{vector}
(\cref{sec:stl}): dobrar.

\begin{lstlisting}[style=bt, language=bash,
    caption={Saída de \texttt{src/strings/}: o mesmo laço, com e sem
             \texttt{reserve}.}]
crescendo char a char: 15→30 30→60 60→120 120→240     ← dobra, igual ao vector
200 chars sem reserve → 4 alocações
200 chars com reserve → 1 alocação
\end{lstlisting}

Repare que a sequência começa em 15, e não em zero: a capacidade inicial de uma
\emph{string} vazia já é o \emph{buffer} da SSO, que existe de qualquer maneira
dentro do objeto.

Assim como no \texttt{vector}, \texttt{clear()} zera o tamanho e \textbf{mantém}
a capacidade, e \texttt{shrink\_to\_fit()} é uma requisição não vinculante. O
idioma que funciona sempre é o mesmo: \texttt{std::string\{\}.swap(s);}.

\subsubsection{Concatenação: por que \texttt{+} custa mais}

A expressão \texttt{a + b + c + d} não é uma operação com quatro argumentos: é
\texttt{((a + b) + c) + d}, e cada \texttt{+} produz uma \emph{string}
\textbf{temporária} nova. Numa cadeia com $N$ trechos, isso tende a $N-1$
objetos intermediários.

Na prática o número não explode, e a razão é interessante: a libstdc++
sobrecarrega \texttt{operator+} para o caso em que o operando da esquerda é um
\gls{categoriadevalor} \emph{rvalue} --- isto é, um temporário --- e nesse caso
reaproveita o \emph{buffer} dele em vez de alocar outro. O subprojeto conta as
alocações das quatro formas de montar o mesmo texto e mostra que a diferença
existe, mas é menor do que a análise ingênua prevê.

\begin{regra}[para $N$ trechos, \texttt{reserve} mais \texttt{append}]
Quando o tamanho final é conhecido ou estimável, \texttt{reserve} seguido de
\texttt{append} (ou de \texttt{+=}) faz exatamente uma alocação. Para texto
formatado, \texttt{std::format} e \texttt{fmt::format} já fazem isso
internamente, e \texttt{format\_to} escreve direto num destino existente ---
ver a \cref{sec:format}.
\end{regra}

\subsubsection{O que \classe{std::string} não sabe}

\classe{std::string} é \texttt{basic\_string<char>}: uma sequência de
\textbf{bytes}. Ela não sabe o que é um caractere Unicode, e \texttt{size()}
conta bytes:

\begin{lstlisting}[style=bt, language=bash,
    caption={Saída de \texttt{src/strings/}: bytes não são letras.}]
"cafe".size() = 4     (4 letras, 4 bytes)
"café".size() = 5     (4 letras, 5 BYTES)
\end{lstlisting}

O acento em UTF-8 ocupa dois bytes (\texttt{C3~A9}), e a consequência atinge
todas as operações posicionais: \texttt{s[i]} pode cair no meio de um caractere,
\texttt{substr} pode cortá-lo ao meio, e \texttt{std::reverse} sobre UTF-8
produz lixo. Comparar sem normalização Unicode também dá resultados errados,
porque a mesma letra acentuada tem mais de uma representação válida.

\begin{armadilha}[a biblioteca padrão não resolve Unicode, e finge que sim]
Os tipos \texttt{u8string}, \texttt{u16string} e \texttt{u32string} apenas
marcam a intenção: nenhum deles implementa normalização, segmentação de
\emph{grafemas} ou comparação sensível a idioma. O cabeçalho
\cabecalho{codecvt}, que prometia conversão entre codificações, foi
\textbf{depreciado no C++17} e nunca ganhou substituto. Para Unicode de verdade
a resposta está fora da biblioteca padrão: ICU, utf8cpp ou Boost.Locale.
\end{armadilha}

O ponto de extensão que existe é o \texttt{char\_traits}, o segundo parâmetro de
\texttt{basic\_string}, que define comparação, cópia e busca. Trocá-lo permite,
por exemplo, uma \emph{string} insensível a maiúsculas sem alocar --- mas o
resultado é um \textbf{tipo diferente}, que não converte para
\classe{std::string}. Quase sempre é melhor normalizar os dados na entrada.

\subsubsection{\classe{std::string\_view}: a janela que não possui}
\label{sec:string-view}

\begin{porque}[que problema o \emph{view} resolve]
Uma função que aceita \texttt{const std::string\&} obriga quem tem um
\texttt{const char*} --- um literal, um \emph{buffer} de rede, um trecho de
outra \emph{string} --- a \textbf{construir uma \texttt{std::string} inteira} só
para passar o argumento. É uma alocação e uma cópia por chamada, para uma
função que só vai ler. \classe{std::string\_view} é o par (ponteiro, tamanho)
que descreve caracteres pertencentes a outra pessoa: converte implicitamente de
tudo o que já é contíguo, e não aloca nada.
\end{porque}

O ganho é duplo. Na chamada, \texttt{contar\_vogais(literal)} deixa de alocar.
E nas operações de recorte: \texttt{str.substr(4,5)} constrói uma
\emph{string} nova, enquanto \texttt{view.substr(4,5)} apenas ajusta o ponteiro
e o tamanho --- é $O(1)$ e não toca na \emph{heap}.

O preço é que o \emph{view} não possui nada, e disso decorrem duas famílias de
armadilhas.

\paragraph{A ausência do terminador nulo.} Uma \emph{C-string} é uma sequência
de bytes terminada por um byte zero (\texttt{'\textbackslash 0'}); é assim que
\texttt{printf}, \texttt{fopen} e toda API C descobrem onde o texto acaba. Um
\emph{view} não tem esse byte: ele é um recorte de outra coisa, e o byte
seguinte ao último caractere pertence à origem.

\begin{lstlisting}[style=bt, language=bash,
    caption={Saída de \texttt{src/strings/}: o que uma API C leria.}]
parcial = "abc" (size=3)
o que printf("%s", parcial.data()) leria: 'abcdefghij'    ← foi até o fim da ORIGEM
\end{lstlisting}

No exemplo do subprojeto a origem é uma \classe{std::string}, que por acaso
termina em nulo, e a leitura para em dez caracteres. Se a origem fosse um
\emph{buffer} binário, a leitura seguiria memória adentro até encontrar um zero
qualquer --- leitura fora de limites, que o AddressSanitizer
(\cref{sec:sanitizers}) detectaria.

\begin{table}[htbp]
\centering
\caption{As três formas corretas de levar um \emph{view} a uma API que espera
         C-string.}
\label{tab:view-c-api}
\footnotesize
\begin{tabular}{@{}ll@{}}
\toprule
\textbf{Forma} & \textbf{Como} \\
\midrule
(a) copiar        & \texttt{std::string\{v\}.c\_str()} --- aloca, mas é correto \\
(b) passar o par  & \texttt{printf("\%.*s", (int)v.size(), v.data())} \\
(c) garantir na origem & assegurar que o \emph{view} cobre até um
                         \texttt{'\textbackslash 0'} \\
\bottomrule
\end{tabular}
\end{table}

\begin{regra}[\texttt{.data()} de um \emph{view} nunca vai direto para uma API C]
A assinatura de \texttt{.data()} é a mesma de \texttt{std::string::data()} ---
\texttt{const char*} --- e é justamente por serem indistinguíveis que o erro
passa na revisão. Se a função de destino não recebe também o tamanho, o
\emph{view} precisa virar \classe{std::string} antes.
\end{regra}

\paragraph{Tempo de vida.} O \emph{view} é tão frágil quanto um ponteiro, porque
\emph{é} um ponteiro. Os três erros clássicos são: apontar para um temporário
que morre no fim da expressão; devolver um \emph{view} para uma variável local;
e guardar um \emph{view} como membro sem documentar quem é o dono do texto.

\begin{lstlisting}[style=bt, language=C++,
    caption={O caso seguro em \texttt{src/strings/demo\_string\_view.hpp}: o
             \emph{view} aponta para um membro.}]
class Documento {
public:
    explicit Documento(std::string conteudo) : conteudo_(std::move(conteudo)) {}

    // Seguro: o view aponta para o membro, cuja vida é a do Documento
    std::string_view titulo() const {
        const auto fim = conteudo_.find('\n');
        return std::string_view{conteudo_}.substr(0, fim);
    }

private:
    std::string conteudo_;
};
\end{lstlisting}

\begin{regra}[as três regras de tempo de vida de um \emph{view}]
(1)~\emph{view} é \textbf{parâmetro} --- o uso dominante, e sempre seguro,
porque o argumento vive até o fim da expressão que contém a chamada.
(2)~\emph{view} como \textbf{retorno} só se a origem sobrevive à chamada, como
no \texttt{Documento} acima. (3)~\emph{view} como \textbf{membro} exige
documentar quem é o dono --- e, quase sempre, é melhor guardar uma
\classe{std::string}.
\end{regra}

\paragraph{As conversões, e a que custa.} \texttt{string $\rightarrow$
string\_view} é implícita e barata. \texttt{string\_view $\rightarrow$ string} é
\textbf{explícita}, porque aloca: \texttt{std::string t = v;} não compila, e
\texttt{std::string t\{v\};} compila. O desenho é intencional --- a conversão
cara precisa ser visível no código.

Há, porém, uma conversão implícita que custa: \texttt{const char* $\rightarrow$
string\_view} chama \texttt{strlen}, que é $O(n)$. Num laço quente com
literais, isso se paga a cada iteração. A saída é o literal \texttt{sv} do
C++17, cujo tamanho é \gls{constexpr}:

\begin{lstlisting}[style=bt, language=C++,
    caption={Transcrição de \texttt{src/strings/demo\_string\_view.hpp}: o
             literal cujo tamanho é conhecido em compilação.}]
    using namespace std::string_view_literals;
    constexpr auto literal_sv = "tamanho conhecido na compilação"sv;

    // ⚠ 33, não 31: a frase tem 31 CARACTERES, mas o 'ç' e o 'ã' ocupam 2
    //   bytes cada em UTF-8. `size()` conta bytes — é a mesma lição de
    //   demo_string.hpp, seção 1.5, e aqui ela foi pega por um static_assert.
    static_assert(literal_sv.size() == 33);
\end{lstlisting}

O comentário registra um caso em que a medição contrariou a expectativa do
próprio autor: o \texttt{static\_assert} foi escrito com 31, o valor que a
contagem de letras sugere, e o compilador recusou. São 33 bytes, porque
\texttt{ç} e \texttt{ã} ocupam dois cada --- exatamente a lição da subseção
anterior, aplicada a quem já a tinha escrito.

\begin{table}[htbp]
\centering
\caption{Escolhendo a assinatura de um parâmetro de texto.}
\label{tab:assinatura-texto}
\footnotesize
\begin{tabular}{@{}ll@{}}
\toprule
\textbf{Assinatura} & \textbf{Quando} \\
\midrule
\texttt{std::string\_view}   & só \textbf{lê}, e não guarda --- o padrão \\
\texttt{const std::string\&} & precisa de \texttt{.c\_str()} garantido \\
\texttt{std::string} por valor & vai \textbf{guardar} uma cópia (mais \texttt{std::move}) \\
\texttt{std::string\&}       & vai \textbf{modificar} a do chamador \\
\bottomrule
\end{tabular}
\end{table}

\subsubsection{\texttt{std::format}: formatação verificada em compilação}
\label{sec:format}

\begin{porque}[por que uma terceira forma de formatar]
O C++ tinha duas, e as duas eram ruins. O \texttt{printf} é conciso e rápido, e
não tem segurança de tipo: \texttt{printf("\%d", 1.5)} compila e imprime lixo,
\texttt{printf("\%s", std::string\{\})} compila e é \gls{ub}, e
\texttt{printf("\%d \%d", 1)} lê a pilha. Os \emph{iostreams} são seguros e
verbosíssimos, e seus manipuladores são \textbf{estado global e pegajoso}: um
\texttt{std::hex} esquecido contamina toda a saída seguinte. \texttt{std::format}
(C++20) tem a sintaxe do primeiro, a segurança do segundo e uma propriedade que
nenhum dos dois tem --- a \emph{string} de formato é validada na
\textbf{compilação}.
\end{porque}

O subprojeto usa \texttt{fmt::format}, a biblioteca que originou a proposta,
porque é o que já está no \texttt{conanfile.py} do repositório
(\cref{sec:ambiente}); a API é a mesma.

\paragraph{O que a validação pega.} As três linhas abaixo não compilam:

\begin{lstlisting}[style=bt, language=C++,
    caption={Erros de formato que viram erro de \emph{build}.}]
fmt::format("{:d}", "texto")   // NÃO COMPILA — 'd' não se aplica a texto
fmt::format("{} {}", 1)        // NÃO COMPILA — dois campos, um argumento
fmt::format("{", 1)            // NÃO COMPILA — chave não fechada
\end{lstlisting}

O mecanismo é o \texttt{consteval} (\cref{sec:compile-time}): a \emph{string} de
formato é analisada em tempo de compilação, e um \texttt{throw} durante uma
avaliação constante é, por definição, erro de \emph{build}. Não há custo em
tempo de execução --- a análise já aconteceu.

\begin{armadilha}[a validação pega o autor deste livro também]
O demo original continha um \texttt{\{:.1\%\}} para imprimir uma porcentagem.
O \texttt{\%} não é especificador válido para \texttt{double} nem na fmt~10 nem
no \texttt{std::format} do C++20, e o \textbf{arquivo não compilou}. É o
comportamento anunciado funcionando contra quem o anunciou: com
\texttt{printf}, o mesmo erro teria produzido saída errada em silêncio. A
porcentagem se faz multiplicando por 100 antes e usando \texttt{\{:.1f\}\%}.
\end{armadilha}

\paragraph{A mini-linguagem.} A forma geral de um campo de substituição é
\texttt{\{índice:preenchimento alinhamento sinal \# 0 largura .precisão tipo\}}.
Vale memorizar os grupos, não a gramática:

\begin{lstlisting}[style=bt, language=bash,
    caption={Saída de \texttt{src/strings/}: os grupos da especificação.}]
{:<10}  {:>10}  {:^10}  {:*^10}      alinhamento e preenchimento
{:08.3f}  {:+}  {:.2e}  {:g}          números
{:b}  {:#b}  {:o}  {:#x}  {:#X}       bases
{0} {1} {0}                            índices explícitos e reuso
{:{}.{}f}                              largura/precisão dinâmicas
\end{lstlisting}

Duas capacidades merecem nota porque não têm equivalente em \texttt{printf}: o
\textbf{reuso de argumento} por índice explícito (\texttt{\{0\} \{1\} \{0\}}
imprime o primeiro argumento duas vezes, sem passá-lo duas vezes) e a
formatação de tipos da biblioteca padrão --- \texttt{vector}, \texttt{pair},
\texttt{map} --- que na fmt vem de \cabecalho{fmt/ranges.h} e no padrão chegou
no C++23.

\paragraph{Formatar sem alocar.} \texttt{format} devolve uma
\classe{std::string} nova, portanto aloca. Quando isso importa, a família tem
variantes que escrevem num destino existente:

\begin{table}[htbp]
\centering
\caption{As quatro portas de saída da biblioteca de formatação.}
\label{tab:format-destinos}
\footnotesize
\begin{tabular}{@{}lll@{}}
\toprule
\textbf{Função} & \textbf{Destino} & \textbf{Aloca} \\
\midrule
\texttt{format}      & \classe{std::string} nova  & sim \\
\texttt{format\_to}  & qualquer iterador de saída & não$^*$ \\
\texttt{format\_to\_n} & idem, com limite de bytes & não \\
\texttt{print}       & \emph{stdout} direto        & não \\
\bottomrule
\end{tabular}\\[2pt]
\footnotesize $^*$ depende do destino: um \texttt{back\_inserter} sobre um
\texttt{vector} pode realocar.
\end{table}

\texttt{format\_to\_n} escrevendo num vetor de \texttt{char} na pilha é o
equivalente seguro do \texttt{snprintf}; \texttt{fmt::print} escreve direto no
\emph{stdout} sem construir a \emph{string} intermediária, e é por isso que ele
aparece em todo o repositório em vez de
\texttt{std::cout << fmt::format(...)}.

\paragraph{Um \texttt{formatter} para o seu tipo.} Ensinar a biblioteca a
formatar um tipo próprio é especializar \texttt{std::formatter<T>} (ou
\texttt{fmt::formatter<T>}) com dois métodos: \texttt{parse}, que consome os
especificadores entre \texttt{:} e \texttt{\}}, e \texttt{format}, que escreve
no iterador de saída.

\begin{lstlisting}[style=bt, language=C++,
    caption={Transcrição de \texttt{src/strings/demo\_format.hpp}: um
             \emph{formatter} com especificadores próprios.}]
template <>
struct fmt::formatter<demo_format::Temperatura> {
    char unidade = 'C';

    constexpr auto parse(format_parse_context& ctx) {
        auto it = ctx.begin();
        if (it != ctx.end() && (*it == 'C' || *it == 'F' || *it == 'K')) {
            unidade = *it++;
        }
        if (it != ctx.end() && *it != '}') {
            throw format_error("Temperatura aceita apenas C, F ou K");
        }
        return it;
    }

    auto format(const demo_format::Temperatura& t, format_context& ctx) const {
        switch (unidade) {
            case 'F': return fmt::format_to(ctx.out(), "{:.1f}°F", t.celsius * 9 / 5 + 32);
            case 'K': return fmt::format_to(ctx.out(), "{:.2f}K",  t.celsius + 273.15);
            default:  return fmt::format_to(ctx.out(), "{:.1f}°C", t.celsius);
        }
    }
};
\end{lstlisting}

O detalhe decisivo é o \texttt{constexpr} no \texttt{parse}: como a análise
acontece em tempo de compilação, o \texttt{throw} daquele terceiro \texttt{if}
não é uma exceção em tempo de execução --- é um \textbf{erro de compilação}.
\texttt{fmt::print("\{:F\}", t)} imprime \texttt{77.0°F}; \texttt{"\{:X\}"} não
compila. O seu tipo ganha a mesma verificação que os tipos nativos têm.

O atalho, quando o tipo é essencialmente um número, é herdar de um
\emph{formatter} existente: \texttt{struct formatter<Metros> :
formatter<double>} reaproveita toda a mini-linguagem, e \texttt{\{:8.2f\}}
passa a funcionar de graça.

\paragraph{O custo, e a lição que vale mais que ele.} O subprojeto cronometra
quatro formas de produzir a mesma linha, cem mil vezes --- e o resultado
\textbf{não} mostra a fmt ganhando:

\begin{lstlisting}[style=bt, language=bash,
    caption={Saída de \texttt{src/strings/}: 100\,000 formatações, nas duas
             configurações de compilação.}]
-O0: snprintf 18ms · format_to_n 78ms · format 91ms
-O2: snprintf 18ms · format_to_n 53ms · format 57ms
\end{lstlisting}

O demo explica por quê, em vez de esconder: este repositório linka
\texttt{libfmtd}, a variante de \textbf{depuração} que o Conan compilou com
\texttt{build\_type=Debug}. Compilar o arquivo do demo com \texttt{-O2} quase
não muda o quadro, porque o custo está dentro da biblioteca --- que já foi
compilada. O \texttt{snprintf}, por sua vez, vem da glibc, sempre otimizada.

\begin{regra}[um número sem as \emph{flags} ao lado é um boato]
Três exigências decorrem do caso acima, e nenhuma delas é sobre a fmt.
(1)~Saiba em que modo a \textbf{biblioteca} foi compilada, não só o seu código.
(2)~Compare pares que diferem em \textbf{uma} variável só --- aqui, a única
comparação controlada é \texttt{format\_to\_n} contra \texttt{format}, que
diferem apenas no destino e isolam o custo de alocar uma \emph{string} por
iteração. (3)~Desconfie de qualquer \emph{benchmark} cujo autor não declara as
\emph{flags} --- inclusive deste.
\end{regra}

E há o argumento que não depende de medição alguma: \texttt{snprintf} não tem
segurança de tipo, e o \texttt{ostringstream} carrega estado global pegajoso.
\texttt{format} dá as duas coisas.

\subsubsection{\cabecalho{charconv}: texto e número sem alocar e sem locale}
\label{sec:charconv}

Converter \texttt{"42"} em \texttt{42} parece trivial, e o C++ acumulou cinco
formas de fazê-lo --- todas com algum defeito.

\begin{table}[htbp]
\centering
\caption{As cinco formas de converter texto em número, e o defeito de cada uma.}
\label{tab:texto-numero}
\footnotesize
\begin{tabular}{@{}ll@{}}
\toprule
\textbf{Forma} & \textbf{Defeito} \\
\midrule
\texttt{atoi}         & sem detecção de erro --- \texttt{"abc"} devolve 0, igual a \texttt{"0"} \\
\texttt{strtol}       & detecta, via a variável global \texttt{errno} \\
\texttt{stoi}         & \textbf{lança} --- caro no caminho de erro \\
\texttt{stringstream} & aloca, e respeita o \emph{locale} \\
\texttt{from\_chars}  & --- \\
\bottomrule
\end{tabular}
\end{table}

\emph{Locale} é a configuração global de convenções culturais do processo:
separador decimal, separador de milhar, nomes de meses. É definido por
\texttt{setlocale} e afeta \texttt{printf}, os \emph{iostreams} e as funções
\texttt{sto*} --- inclusive quando quem o mudou foi uma biblioteca de terceiros
na inicialização.

\paragraph{\texttt{from\_chars}.} A assinatura recebe um intervalo
\texttt{[primeiro, último)} de \texttt{char} e uma referência ao destino, e
devolve uma \emph{struct} com dois campos: \texttt{ptr}, onde a conversão parou,
e \texttt{ec}, um \texttt{std::errc}. Os casos:

\begin{lstlisting}[style=bt, language=bash,
    caption={Saída de \texttt{src/strings/}: \texttt{from\_chars} sobre entradas
             típicas.}]
'42'             → 42 ✓
'3.14'           → 3 — mas SOBROU '.14'      ← o ptr permite detectar
'42abc'          → 42 — mas SOBROU 'abc'      ← stoi devolveria 42 calado
' 42'            → invalid_argument            ← não pula espaços
'99999999999999' → result_out_of_range
\end{lstlisting}

O \texttt{ptr} devolvido é a diferença prática mais importante:
\texttt{stoi("42abc")} devolve 42 sem reclamar, e não há como o chamador
descobrir que sobrou lixo. Com \texttt{from\_chars}, comparar \texttt{ptr} com o
fim do intervalo responde exatamente isso. E é o mesmo \texttt{ptr} que
realimenta o laço quando se \emph{parseia} uma sequência inteira --- o padrão
que sustenta os \emph{parsers} rápidos de JSON e CSV.

A lista do que \texttt{from\_chars} \textbf{não} faz é a lista de suas
qualidades: não pula espaços, não aceita o prefixo \texttt{0x} (a base vai por
parâmetro), não aceita \texttt{+} inicial, não aloca, não lança e não olha o
\emph{locale}. Cada ``não'' é uma política devolvida a quem chama, que passa a
pagar por ela só quando precisa.

\paragraph{\texttt{to\_chars} e o \emph{round-trip} garantido.} Na direção
inversa há uma garantia que nenhuma alternativa oferece: chamado \textbf{sem}
especificar precisão, \texttt{to\_chars} produz a \textbf{menor} \emph{string}
que faz \texttt{from\_chars} devolver exatamente o mesmo \texttt{double}.

\begin{lstlisting}[style=bt, language=bash,
    caption={Saída de \texttt{src/strings/}: ida e volta exata.}]
0.1                → '0.1'                ✓ round-trip exato
0.3333333333333333 → '0.3333333333333333' ✓
123456789.12345679 → '123456789.12345679' ✓
\end{lstlisting}

O algoritmo por trás disso (Ryū / Grisu) resolve um problema real:
\texttt{printf("\%g")} usa seis dígitos significativos e perde precisão, e
\texttt{printf("\%.17g")} sempre imprime dezessete, produzindo
\texttt{0.10000000000000001} para \texttt{0.1}. Onde o número precisa voltar
idêntico depois de virar texto --- serialização JSON, telemetria, formatos
textuais --- essa é a ferramenta.

\begin{armadilha}[o \emph{locale} corrompe dados em silêncio]
Com o \emph{locale} global em \texttt{pt\_BR} ou \texttt{de\_DE}, o separador
decimal é a vírgula:
\begin{lstlisting}[style=bt, language=bash, frame=none,
                   xleftmargin=0pt, xrightmargin=0pt]
ostringstream << 3.14  → "3,14"     ← VÍRGULA
stod("3.14")           → 3.0        ← para no ponto, e não reclama
\end{lstlisting}
O bug completo tem quatro passos: o programa grava um CSV numa máquina
\texttt{en\_US}, outra máquina em \texttt{pt\_BR} lê o mesmo arquivo,
\texttt{stod("3.14")} devolve 3.0 sem erro, e o dado fica errado para sempre. E
basta uma biblioteca de terceiros chamar \texttt{setlocale(LC\_ALL, "")} para o
comportamento do \emph{seu} código mudar. \cabecalho{charconv} é imune por
projeto: não tem parâmetro de \emph{locale}, e o separador é sempre o ponto.
\end{armadilha}

\begin{regra}[\emph{locale} é apresentação, não dado]
Para \textbf{dados} --- arquivos, rede, protocolos, formatos de intercâmbio ---
use \cabecalho{charconv}. Para \textbf{apresentação} ao usuário, use
\texttt{std::format} com o especificador \texttt{\{:L\}}, que pede a formatação
sensível a \emph{locale} \emph{explicitamente}, no ponto em que ela é desejada.
\end{regra}

O subprojeto também cronometra as três formas convertendo a mesma
\classe{std::string} no mesmo \texttt{int}, e aqui a comparação \emph{é}
controlada. As diferenças são estruturais e não dependem de \emph{build}:
\texttt{stoi} constrói uma \emph{string} temporária internamente e usa
\texttt{strtol} com \texttt{errno}; \texttt{istringstream} constrói um objeto de
\emph{stream} inteiro por conversão, com \emph{locale}, \emph{sentry} e
\emph{flags}; \texttt{from\_chars} é um laço sobre os dígitos. Há ainda um custo
que o cronômetro não vê: \texttt{stoi} \textbf{lança} em entrada inválida, e num
\emph{parser} que encontra lixo com frequência cada erro custa um
desempilhamento completo --- medido em \cref{sec:errors} como cerca de mil vezes
mais caro que o caminho feliz.

\subsubsection{\classe{std::regex}: o custo real}
\label{sec:regex}

Uma expressão regular é uma notação para descrever conjuntos de cadeias de
caracteres. \classe{std::regex} chegou no C++11 e é, com folga, a parte mais
criticada da biblioteca padrão: a API é razoável, a implementação é
notoriamente lenta.

\paragraph{\texttt{match} contra \texttt{search}.} É o erro número um, e vem de
quem conhece outras linguagens:

\begin{table}[htbp]
\centering
\caption{As três formas de aplicar uma expressão regular.}
\label{tab:regex-funcoes}
\footnotesize
\begin{tabular}{@{}ll@{}}
\toprule
\textbf{Função} & \textbf{Casa} \\
\midrule
\texttt{regex\_match}     & a \emph{string} \textbf{inteira} \\
\texttt{regex\_search}    & a primeira ocorrência, em qualquer posição \\
\texttt{sregex\_iterator} & todas as ocorrências \\
\bottomrule
\end{tabular}
\end{table}

Em Python, \texttt{re.match} casa no \emph{início} da \emph{string}; em C++,
\texttt{regex\_match} exige a \emph{string} toda. Sobre
\texttt{"pedido 42 entregue"} com o padrão \texttt{\textbackslash d+},
\texttt{regex\_match} devolve falso e \texttt{regex\_search} devolve verdadeiro.

\paragraph{Capturas.} Grupos entre parênteses são recuperados por um objeto
\texttt{std::smatch}: \texttt{m[0]} é o casamento inteiro, \texttt{m[1]} em
diante são os grupos, e \texttt{m.prefix()} e \texttt{m.suffix()} são o texto
antes e depois. Em \texttt{regex\_replace}, o formato de substituição usa
\texttt{\$1} a \texttt{\$9} para os grupos, \texttt{\$\&} para o casamento
inteiro e \texttt{\$\$} para um cifrão literal.

\begin{armadilha}[\texttt{smatch} guarda iteradores, não cópias]
Um \texttt{std::smatch} não copia o texto: ele guarda \textbf{iteradores} para a
\emph{string} original. Se ela morrer, o \texttt{smatch} fica pendurado, e o
caso mais comum é fatal:
\begin{lstlisting}[style=bt, language=C++, frame=none,
                   xleftmargin=0pt, xrightmargin=0pt]
std::smatch m;
std::regex_search(obterTexto(), m, re);   // ✗ argumento temporário
m[1].str();                               // use-after-free
\end{lstlisting}
O comitê considerou o caso grave o bastante para \textbf{deletar} a sobrecarga
que recebe \texttt{const std::string\&\&}: a primeira linha nem compila. É um
dos poucos lugares da biblioteca padrão em que um erro de tempo de vida virou
erro de compilação por decreto.
\end{armadilha}

\paragraph{Duas regras de escrita.} Sempre use \emph{raw string} para o padrão
--- \texttt{R"(\textbackslash d+)"} em vez de
\texttt{"\textbackslash\textbackslash d+"} ---, porque sem ela o compilador
processa os escapes antes de a expressão chegar à biblioteca. E nunca construa
uma \texttt{std::regex} dentro de um laço: o construtor \emph{compila} a
expressão, e isso é caro.

\paragraph{O custo, medido.} O subprojeto extrai o campo
\texttt{codigo=(\textbackslash d+)} de uma linha de \emph{log} de três formas:

\begin{lstlisting}[style=bt, language=bash,
    caption={Saída de \texttt{src/strings/}: o mesmo campo, três formas.}]
find + from_chars    :   0 ms   (20000 iterações)
regex fora do laço   : 133 ms   (20000 iterações)
regex DENTRO do laço : 176 ms   ( 2000 iterações — 10× MENOS trabalho!)
\end{lstlisting}

São duas lições distintas, e a segunda é a que surpreende. A primeira: construir
a expressão dentro do laço domina tudo --- a terceira linha fez \textbf{dez
vezes menos trabalho} e ainda assim levou mais tempo que a segunda. A segunda:
mesmo pré-construída, para um padrão \textbf{fixo e simples} a expressão regular
não se paga; um \texttt{find} seguido de \texttt{from\_chars} resolve o mesmo
problema num tempo que o cronômetro em milissegundos nem registra.

\begin{table}[htbp]
\centering
\caption{Quando a expressão regular é a ferramenta certa.}
\label{tab:regex-quando}
\footnotesize
\begin{tabular}{@{}ll@{}}
\toprule
\textbf{Use \texttt{std::regex} quando} & \textbf{Não use quando} \\
\midrule
o padrão vem do \textbf{usuário}      & laço quente com padrão fixo \\
a expressão é complexa e roda uma vez & o formato é conhecido: escreva o \emph{parser} \\
legibilidade importa mais que velocidade & precisa de desempenho: RE2, PCRE2, CTRE \\
\bottomrule
\end{tabular}
\end{table}

\begin{armadilha}[ReDoS: entrada de terceiros com padrão de terceiros]
\classe{std::regex} usa \emph{backtracking} --- ao falhar, volta atrás e tenta
outra combinação. Para padrões com quantificadores aninhados, como
\texttt{(a+)+\$}, o número de combinações cresce \textbf{exponencialmente} com o
tamanho da entrada, e uma cadeia de algumas dezenas de caracteres trava o
processo. É o ataque conhecido como ReDoS. A regra é dupla: nunca aplique uma
expressão vinda de terceiros a uma entrada vinda de terceiros; e, se o padrão é
configurável pelo usuário, prefira RE2, que garante tempo linear por construção.
\end{armadilha}

% =============================================================================
\subsection{Tempo}
\label{sec:chrono}
% =============================================================================

\begin{porque}[por que \cabecalho{chrono} não é ``pegar a hora'']
A API C de tempo é um conjunto de inteiros sem unidade: \texttt{time\_t} conta
segundos, \texttt{clock\_t} conta tiques de origem indefinida, e
\texttt{sleep(1)} dorme um segundo enquanto \texttt{usleep(1)} dorme um
microssegundo. Nada no tipo distingue os dois, e passar um pelo outro é um bug
de fator um milhão que compila. \cabecalho{chrono} é a resposta: um
\textbf{sistema de unidades com verificação em compilação}, aplicado ao tempo.
Somar horas com milissegundos funciona e dá o resultado certo; somar horas com
o número 5 não compila, porque 5 não é uma quantidade de tempo.
\end{porque}

\subsubsection{\texttt{duration}: a unidade vive no tipo}

A declaração é:

\begin{lstlisting}[style=bt, language=C++,
    caption={A forma geral de uma duração.}]
template <class Rep, class Period = std::ratio<1>>
class duration;
\end{lstlisting}

\texttt{Rep} é o tipo que guarda a contagem --- \texttt{int64\_t},
\texttt{double}; \texttt{Period} é a fração de \textbf{segundo} que cada tique
representa, expressa como um \texttt{std::ratio}, que é um par de inteiros
(numerador e denominador) codificado no tipo. Assim,
\texttt{std::chrono::milliseconds} é \texttt{duration<int64\_t, ratio<1,1000>>}:
o ``milli'' está no \textbf{tipo}, não em uma variável.

\begin{lstlisting}[style=bt, language=bash,
    caption={Saída de \texttt{src/chrono/}: as unidades e o que elas custam.}]
nanoseconds    Rep=int64_t  Period=1/1000000000
milliseconds   Rep=int64_t  Period=1/1000
hours          Rep=int64_t  Period=3600/1
sizeof(milliseconds) = 8 bytes      ← é só o int64_t
\end{lstlisting}

Toda a segurança dimensional custa \textbf{zero} em tempo de execução: o objeto
é do tamanho do inteiro que ele guarda, e as conversões viram multiplicações
por constantes que o compilador resolve.

O método \texttt{.count()} extrai o número cru e \textbf{descarta a unidade}. Ele
é a saída do sistema de tipos, e deve aparecer só na fronteira --- ao imprimir,
ao gravar, ao chamar uma API C. Usá-lo no meio de um cálculo devolve o problema
que a biblioteca resolve.

\subsubsection{Conversões: implícitas quando são exatas}

A regra é uma só, e é o coração do desenho: \textbf{a conversão é implícita
quando não pode perder informação}.

\begin{lstlisting}[style=bt, language=bash,
    caption={Saída de \texttt{src/chrono/}: o que o compilador aceita e o que
             recusa.}]
2h → milliseconds     ✓ implícito (exato)
1500ms → seconds      ✗ NÃO COMPILA (truncaria)
duration_cast<seconds>(1500ms) = 1 s     ← trunca em direção a ZERO
\end{lstlisting}

Converter para uma unidade \emph{menor} multiplica por um inteiro e é sempre
exato; converter para uma unidade \emph{maior} pode truncar, e o compilador
exige que a perda seja escrita à mão. Quando truncar não é a política desejada,
existem \texttt{floor}, \texttt{ceil} e \texttt{round}, todos do C++17.

\begin{armadilha}[para negativos, truncar e arredondar para baixo são coisas diferentes]
\texttt{duration\_cast} trunca \textbf{em direção a zero};
\texttt{floor} arredonda \textbf{para baixo}. Para valores positivos os dois
coincidem, e é por isso que a diferença passa despercebida em teste:
\begin{lstlisting}[style=bt, language=bash, frame=none,
                   xleftmargin=0pt, xrightmargin=0pt]
duration_cast<seconds>(-1500ms) = -1 s
floor<seconds>(-1500ms)         = -2 s
\end{lstlisting}
Em cálculo de fuso horário, de compensação de deriva ou de qualquer diferença
que possa ser negativa, escolher o errado produz um segundo de erro que só
aparece em metade dos casos.
\end{armadilha}

\subsubsection{Aritmética e unidades derivadas}

Somar durações de unidades diferentes produz o \emph{tipo comum}: a unidade mais
fina que representa exatamente as duas.

\begin{lstlisting}[style=bt, language=bash,
    caption={Saída de \texttt{src/chrono/}: a álgebra das durações.}]
1h + 30min + 45s = 5445 s      ← tipo comum: seconds
1s + 500ms       = 1500 ms     ← tipo comum: milliseconds
1h / 15min       = 4           ← escalar puro, sem unidade
1h * 1h          ✗ tempo² não é um duration
1h + 5           ✗ 5 do quê?
\end{lstlisting}

A terceira linha é a mais instrutiva: dividir tempo por tempo produz um número
\textbf{sem unidade}, e o tipo do resultado é um inteiro, não uma
\texttt{duration}. É a álgebra dimensional correta, obtida de graça pelo sistema
de tipos.

\texttt{Rep} não precisa ser inteiro. Com \texttt{double}, converter para uma
unidade maior deixa de perder informação e passa a ser implícito ---
\texttt{duration<double> s = 1500ms;} compila. A troca é a exatidão: para
contabilidade de tempo (agendamento, \emph{timeouts}), prefira inteiros; para
\emph{relatar} medições, \texttt{double} é mais legível.

Definir uma unidade nova é uma linha, e ela participa de toda a aritmética e de
todas as verificações automaticamente:

\begin{lstlisting}[style=bt, language=C++,
    caption={Transcrição de \texttt{src/chrono/demo\_duration.hpp}: quadros de
             vídeo como unidade de tempo.}]
    using Frame60 = ch::duration<int64_t, std::ratio<1, 60>>;
    using Frame24 = ch::duration<int64_t, std::ratio<1, 24>>;

    const Frame60 um_segundo_de_video = 1s;
\end{lstlisting}

\texttt{1s} converte implicitamente para 60 quadros, porque a conversão é
exata; converter de 60~fps para 24~fps exige \texttt{duration\_cast}, porque
trunca. O compilador aplica a mesma regra a uma unidade que ele nunca viu.

\subsubsection{Os relógios, e por que escolher o errado é um bug}
\label{sec:clocks}

A diferença entre os relógios da biblioteca não é de precisão --- é de
\textbf{garantia}. Um relógio \emph{monotônico} nunca anda para trás; um relógio
de \emph{parede} corresponde à hora do mundo e, por isso mesmo, pode ser
ajustado --- pelo NTP, pelo usuário, pelo horário de verão.

\begin{table}[htbp]
\centering
\caption{Os três relógios clássicos e para que serve cada um.}
\label{tab:relogios}
\footnotesize
\begin{tabular}{@{}llll@{}}
\toprule
\textbf{Relógio} & \textbf{Monotônico} & \textbf{É hora real} & \textbf{Usar para} \\
\midrule
\texttt{system\_clock}          & \textbf{não}    & sim     & \textbf{datar} \\
\texttt{steady\_clock}          & \textbf{sim}    & não     & \textbf{medir} \\
\texttt{high\_resolution\_clock}& depende!        & depende!& \textbf{nada} \\
\bottomrule
\end{tabular}
\end{table}

O bug que a distinção previne é este: com \texttt{system\_clock}, um ajuste do
NTP entre as duas leituras faz a diferença ser \textbf{negativa}. Com
\texttt{steady\_clock}, isso é impossível por construção.

\begin{armadilha}[\texttt{high\_resolution\_clock} é um apelido, e neste ambiente é o errado]
O padrão permite que \texttt{high\_resolution\_clock} seja um \emph{alias} para
outro relógio, e não diz qual. O subprojeto mede:
\begin{lstlisting}[style=bt, language=bash, frame=none,
                   xleftmargin=0pt, xrightmargin=0pt]
é o mesmo que system_clock? true         ← libstdc++
high_resolution_clock::is_steady = false
\end{lstlisting}
Ou seja: na libstdc++, o relógio cujo nome sugere ``o melhor para medir'' é
justamente o \textbf{não monotônico}. Todo código que mede duração com ele neste
compilador está sujeito a resultado negativo. E o mapeamento muda: libstdc++
aponta para \texttt{system\_clock}, libc++ e MSVC apontam para
\texttt{steady\_clock}. O mesmo código muda de comportamento ao trocar de
compilador, e o bug não aparece em teste --- só quando o relógio do sistema é
ajustado. A recomendação, que é também a do cppreference, é não usá-lo nunca.
\end{armadilha}

\paragraph{Resolução do tipo contra resolução real.} O \texttt{period} de um
relógio diz a resolução que o \textbf{tipo} é capaz de representar. A resolução
\emph{efetiva} --- o menor intervalo que duas leituras consecutivas conseguem
distinguir --- é outra coisa, e se mede:

\begin{lstlisting}[style=bt, language=bash,
    caption={Saída de \texttt{src/chrono/}: o tipo promete nanossegundos; a
             realidade é outra.}]
system_clock   tipo: 1/1000000000 s  |  resolução medida: 10 ns
custo de steady_clock::now(): ~19 ns por chamada
\end{lstlisting}

Chamar \texttt{now()} custa cerca de dezenove nanossegundos. Medir algo que dura
menos que isso é \textbf{medir o relógio} --- e é a razão de existir a
\cref{sec:medicao}.

\subsubsection{\texttt{time\_point}: instantes, e a álgebra que faz sentido}

A distinção que organiza a biblioteca inteira é entre \emph{quanto tempo}
(\texttt{duration}) e \emph{quando} (\texttt{time\_point}). Um
\texttt{time\_point} é, por dentro, apenas uma \texttt{duration} contada desde a
\textbf{\emph{epoch}} do seu relógio --- o instante de referência a partir do
qual ele conta. Para \texttt{system\_clock}, o C++20 garante que a \emph{epoch}
é 1970-01-01 00:00:00 UTC; para \texttt{steady\_clock}, é um ponto arbitrário,
tipicamente o momento em que a máquina ligou.

\begin{table}[htbp]
\centering
\caption{A álgebra de instantes e durações.}
\label{tab:algebra-tempo}
\footnotesize
\begin{tabular}{@{}ll@{}}
\toprule
\textbf{Expressão} & \textbf{Resultado} \\
\midrule
\texttt{time\_point - time\_point} & \texttt{duration} --- ``quanto durou'' \\
\texttt{time\_point + duration}    & \texttt{time\_point} --- ``daqui a'' \\
\texttt{duration $\pm$ duration}   & \texttt{duration} \\
\texttt{time\_point + time\_point} & \textbf{erro de compilação} \\
\bottomrule
\end{tabular}
\end{table}

A última linha é a que importa: somar dois instantes não tem significado, e o
sistema de tipos impede. Pela mesma razão, o relógio faz parte do \emph{tipo} do
\texttt{time\_point}, e \texttt{system\_clock::now() - steady\_clock::now()}
não compila --- as \emph{epochs} são diferentes, e a subtração não teria
sentido. A consequência prática é boa: não há como, por acidente, medir duração
com o relógio de parede e datar um evento com o monotônico.

\paragraph{Instante em texto.} A via C ainda aparece em código legado:
\texttt{to\_time\_t}, \texttt{localtime\_r}, \texttt{strftime}. Vale registrar
que \texttt{localtime} \textbf{sem} o sufixo \texttt{\_r} não é seguro em
programas com várias \emph{threads}: ele devolve um ponteiro para um
\emph{buffer} estático compartilhado. A via C++20 é formatar o
\texttt{time\_point} diretamente, com \texttt{\{:\%Y-\%m-\%d\}} e afins; e
\texttt{time\_point\_cast<seconds>} trunca a precisão quando se quer comparar
carimbos sem os nanossegundos.

\begin{regra}[\emph{deadline} em vez de \emph{timeout} sempre que puder aninhar]
Uma API que espera pode receber ``quanto tempo esperar'' (\emph{timeout}) ou
``até quando esperar'' (\emph{deadline}). A diferença aparece em cadeia:
\begin{lstlisting}[style=bt, language=bash, frame=none,
                   xleftmargin=0pt, xrightmargin=0pt]
conectar(timeout=100ms) → autenticar(timeout=100ms) → consultar(timeout=100ms)
Total possível: 300ms, embora o chamador tenha pedido 100.

auto limite = steady_clock::now() + 100ms;
conectar(limite) → autenticar(limite) → consultar(limite)
Total garantido: 100ms.
\end{lstlisting}
É por isso que a biblioteca oferece os dois: \texttt{wait\_for} recebe duração,
\texttt{wait\_until} recebe instante (\cref{sec:concurrency}). Use
\texttt{\_until} em qualquer função que possa ser chamada dentro de outra, e
\texttt{\_for} apenas na folha da cadeia.
\end{regra}

\subsubsection{O calendário do C++20}

Até o C++17, trabalhar com datas significava \texttt{std::tm}: uma \emph{struct}
C de campos \texttt{int} em que o mês vai de 0 a 11, o ano é ``anos desde
1900'', e nada impede escrever \texttt{tm.tm\_mon = 47}. O C++20 trouxe um
sistema de tipos completo: \texttt{year}, \texttt{month}, \texttt{day},
\texttt{weekday}, e as combinações entre eles.

A sintaxe de composição usa \texttt{operator/}, escolhido por já lembrar data em
quase toda notação do mundo, e aceita as três ordens culturais:

\begin{lstlisting}[style=bt, language=C++,
    caption={As três ordens, todas válidas e todas o mesmo dia.}]
2026y / August / 25d
August / 25d / 2026y
25d / August / 2026y
\end{lstlisting}

Os tipos são minúsculos --- a segurança não custa memória --- e trazem
validação que \texttt{std::tm} nunca teve: \texttt{ok()} responde se a data
existe, e \texttt{year::is\_leap()} implementa a regra bissexta completa
(divisível por 4, exceto os divisíveis por 100 que não sejam por 400 --- por
isso 1900 e 2100 não são bissextos, e 2000 é).

\begin{armadilha}[\texttt{ok()} não é chamado automaticamente]
Construir uma data inválida \textbf{compila e não lança}:
\texttt{2026y/February/30d} é um objeto perfeitamente formado cujo
\texttt{ok()} devolve falso. A biblioteca deixa a verificação para você, e o
lugar de fazê-la é a fronteira de entrada --- ao ler de um arquivo, de um
formulário, de uma requisição.
\end{armadilha}

\paragraph{\texttt{sys\_days}: a ponte.} \texttt{year\_month\_day} é o tipo do
\emph{calendário}: mês, ano, validação. \texttt{sys\_days} é um
\texttt{time\_point<system\_clock, days>}, e é o tipo da \emph{aritmética}:
somar dias, calcular diferenças, descobrir o dia da semana. Converter entre os
dois é barato, e o padrão de uso é converter conforme a operação.

\begin{armadilha}[somar meses pode produzir uma data que não existe]
Somar dias nunca produz data inválida; somar meses ou anos, sim:
\begin{lstlisting}[style=bt, language=bash, frame=none,
                   xleftmargin=0pt, xrightmargin=0pt]
2026-01-31 + 1 mês   = 2026-02-31 is not a valid date  → ok() = false
2026-01-31 + 31 dias = 2026-03-03                       → ok() = true
\end{lstlisting}
A biblioteca \textbf{não conserta}: ela devolve a data inválida e deixa a
política com você. As três políticas possíveis são ir para o último dia do mês
(28/02), transbordar --- que é o que \texttt{sys\_days\{data\_inválida\}} faz,
normalizando para 03/03 --- ou rejeitar, porque \texttt{ok()} é falso. O mesmo
vale para aniversários de quem nasceu em 29 de fevereiro:
\texttt{2024-02-29 + 1 ano} é inválido.
\end{armadilha}

\paragraph{Fusos horários.} A API existe --- \texttt{current\_zone},
\texttt{zoned\_time} --- mas depende do banco \texttt{tzdata} e do suporte
compilado na libstdc++; no GCC~13 ela pode lançar \texttt{std::runtime\_error}
ou falhar na ligação, e por isso o demo não a chama. A regra que não depende de
suporte nenhum é outra, e é a mais importante da subseção:

\begin{regra}[armazene em UTC, converta para local só na apresentação]
Guardar hora local em banco de dados é a origem de metade dos \emph{bugs} de
horário de verão que existem: uma hora do ano acontece duas vezes e outra não
acontece, e não há como distinguir os dois carimbos depois de gravados.
Armazene \texttt{system\_clock} (que é UTC), e converta apenas ao exibir.
\end{regra}

\subsubsection{Como cronometrar sem se enganar}
\label{sec:medicao}

Esta é a subseção que sustenta todos os números do livro. São cinco armadilhas,
em ordem de frequência.

\paragraph{1. O compilador elimina o que você quer medir.} Chama-se
\emph{eliminação de código morto}: se o resultado de um cálculo não é usado, o
compilador tem o direito de apagar o cálculo. Em \texttt{-O2}, um laço que soma
dois milhões de números e descarta o resultado simplesmente \textbf{não existe}
no binário, e o cronômetro mede zero.

\begin{table}[htbp]
\centering
\caption{As três defesas contra a eliminação, em ordem de preferência.}
\label{tab:defesas-dce}
\footnotesize
\begin{tabular}{@{}ll@{}}
\toprule
\textbf{Defesa} & \textbf{Qualidade} \\
\midrule
barreira de otimização (\texttt{asm volatile}) & \textbf{a correta} \\
\texttt{volatile} na variável de resultado     & simples, mas impede otimizações legítimas em volta \\
imprimir ou somar o resultado                  & polui a medida \\
\bottomrule
\end{tabular}
\end{table}

A barreira cabe em três linhas e é o que o Google Benchmark chama de
\texttt{DoNotOptimize}:

\begin{lstlisting}[style=bt, language=C++,
    caption={Transcrição de \texttt{src/chrono/demo\_medicao.hpp}: a barreira de
             otimização.}]
template <typename T>
inline void nao_otimize(T&& valor) {
    // Diz ao compilador: "este valor pode ser lido e modificado por código
    // que você não vê". Ele então precisa materializá-lo de verdade.
    asm volatile("" : : "r,m"(valor) : "memory");
}
\end{lstlisting}

\paragraph{2. O trecho é curto demais.} Como \texttt{now()} custa cerca de vinte
nanossegundos, medir uma operação de poucos nanossegundos mede o relógio:

\begin{lstlisting}[style=bt, language=bash,
    caption={Saída de \texttt{src/chrono/}: a mesma operação, medida de duas
             formas.}]
uma execução medida direto : ~20 ns    ← mediu o RELÓGIO
1000000 execuções / N      :   ~2 ns   ← mediu a operação
\end{lstlisting}

A regra prática: o trecho medido precisa durar pelo menos cem vezes a resolução
do relógio. Com vinte nanossegundos de sobrecarga, isso significa alvos de dois
microssegundos ou mais --- ou repetir e dividir.

\paragraph{3. Aquecimento.} A primeira execução paga o que as seguintes não
pagam: \emph{page faults} --- a materialização, pelo núcleo do sistema
operacional, de páginas de memória que foram reservadas mas nunca tocadas ---,
cache frio, preditor de saltos sem histórico, CPU ainda em frequência baixa.

O demo mede dois casos, e é honesto quando o efeito não aparece: no caso em que
um \texttt{std::iota} preencheu o vetor antes da medição, as cinco execuções
saem parecidas, porque o \texttt{iota} já pagou tudo. A conclusão é a mesma nos
dois casos --- \emph{descarte as primeiras execuções}. O custo de descartar é
zero; o de não descartar é uma conclusão errada.

\paragraph{4. Mediana, não média.}

\begin{lstlisting}[style=bt, language=bash,
    caption={Saída de \texttt{src/chrono/}: 21 execuções do mesmo trabalho.}]
mínimo  : 221 µs   ← o mais próximo do custo REAL
mediana : 225 µs   ← robusta a outliers
média   : 228 µs   ← puxada por qualquer pico
p95     : 240 µs
máximo  : 261 µs   ← o escalonador interferiu
\end{lstlisting}

A média é a pior estatística possível aqui, porque um único pico --- uma troca
de contexto, uma interrupção, outro processo acordando --- a contamina inteira.

\begin{table}[htbp]
\centering
\caption{Qual estatística responde a qual pergunta.}
\label{tab:estatistica-bench}
\footnotesize
\begin{tabular}{@{}ll@{}}
\toprule
\textbf{Você quer} & \textbf{Use} \\
\midrule
o custo típico             & mediana \\
o melhor caso real         & mínimo \\
a latência de um contrato de serviço & p95 / p99 \\
---                        & média (quase nunca) \\
\bottomrule
\end{tabular}
\end{table}

\paragraph{5. Um cronômetro reutilizável.} O subprojeto fecha com um cronômetro
\gls{raii} e uma função de medição que fazem o básico certo:

\begin{lstlisting}[style=bt, language=C++,
    caption={Transcrição de \texttt{src/chrono/demo\_medicao.hpp}: aquecimento,
             repetição e barreira, em oito linhas.}]
template <typename F>
ch::nanoseconds medir(F&& f, int repeticoes = 1) {
    // Aquecimento: descarta a primeira
    { auto r = f(); nao_otimize(r); }

    const auto t0 = ch::steady_clock::now();
    for (int i = 0; i < repeticoes; ++i) { auto r = f(); nao_otimize(r); }
    const auto dt = ch::steady_clock::now() - t0;

    return ch::duration_cast<ch::nanoseconds>(dt) / repeticoes;
}
\end{lstlisting}

E o subprojeto lista o que ele \textbf{não} faz, e que um \emph{benchmark} sério
faria: estatística (mediana, desvio, intervalo de confiança), ajuste automático
de $N$ até a medida estabilizar, fixação da \emph{thread} numa CPU com
\emph{turbo boost} desligado, detecção de interferência do escalonador. Para
isso existem Google Benchmark, nanobench e o suporte do Catch2.

% =============================================================================
\subsection{Bits}
\label{sec:bits}
% =============================================================================

Os operadores \texttt{\&}, \texttt{|}, \texttt{\textasciicircum},
\texttt{\textasciitilde}, \texttt{<<} e \texttt{>>} parecem os mais simples da
linguagem, e escondem três regras que produzem \gls{ub}. O subprojeto
\projeto{bits} abre com as três, resumidas em três linhas:

\begin{lstlisting}[style=bt, language=C++,
    caption={As três armadilhas dos operadores de bit.}]
uint8_t x = 0x0F;
~x                    // NÃO é 0xF0 — é -16, um int    (promoção inteira)
uint64_t m = 1 << 40; // ✗ UB: o literal `1` é int      (shift fora de faixa)
int f : 1; f = 1;     // f vale -1                       (bit de sinal)
\end{lstlisting}

\subsubsection{As quatro operações com máscara}

Uma \emph{máscara} é um valor em que os bits de interesse estão em 1. Com ela,
as quatro operações fundamentais se escrevem sempre da mesma forma:

\begin{table}[htbp]
\centering
\caption{Os quatro idiomas com máscara.}
\label{tab:mascaras}
\footnotesize
\begin{tabular}{@{}ll@{}}
\toprule
\textbf{Operação} & \textbf{Idioma} \\
\midrule
testar   & \texttt{(x \& M) != 0} \\
ligar    & \texttt{x |= M} \\
desligar & \texttt{x \&= \textasciitilde M} \\
alternar & \texttt{x \textasciicircum= M} \\
\bottomrule
\end{tabular}
\end{table}

\subsubsection{A promoção inteira, que morde}

\begin{porque}[por que existe a promoção inteira]
Herdada do C, a regra diz que todo operando de tipo \emph{menor} que
\texttt{int} --- \texttt{char}, \texttt{short}, \texttt{uint8\_t} --- é
convertido para \texttt{int} \textbf{antes} da operação. A motivação era
histórica e razoável: as unidades aritméticas trabalham na largura da palavra, e
uniformizar os operandos simplifica a linguagem. A consequência é que uma
operação escrita sobre oito bits acontece sobre trinta e dois.
\end{porque}

\begin{lstlisting}[style=bt, language=bash,
    caption={Saída de \texttt{src/bits/}: a promoção, medida.}]
uint8_t x = 0x0F;
~x como int       : -16 (0xFFFFFFF0)      ← não é 0xF0
uint8_t(0xFF) << 4 = 0xFF0 (12 bits!)     ← não trunca
\end{lstlisting}

O bug real que isso causa é uma comparação que nunca é verdadeira:
\texttt{if (flags == \textasciitilde MASCARA)} compara um \texttt{uint8\_t} com
um \texttt{int} \textbf{negativo}. O código compila, roda e nunca entra no
\texttt{if}.

As defesas são três: use \texttt{unsigned} ou \texttt{uint32\_t} para máscaras;
faça o \texttt{static\_cast} de volta ao tipo pequeno, explicitamente; e compile
com \texttt{-Wconversion}, para que o compilador cobre a conversão em vez de
fazê-la calado.

\subsubsection{\emph{Shift} fora de faixa}

A regra é: \texttt{x << n} e \texttt{x >> n} são \gls{ub} se \texttt{n} for
negativo ou maior ou igual ao número de bits de \texttt{x} \textbf{após a
promoção}. As consequências são pouco intuitivas:

\begin{lstlisting}[style=bt, language=C++,
    caption={Quatro deslocamentos: dois corretos, dois \gls{ub}.}]
uint32_t x = 1;  x << 32;    // ✗ UB (32 bits, deslocamento de 32)
uint8_t  b = 1;  b << 8;     // ✓ ok! b virou int (32 bits)
uint64_t m = 1 << 40;        // ✗ UB — o literal `1` é int
uint64_t m = 1ull << 40;     // ✓
\end{lstlisting}

A segunda linha é o caso em que a medição contraria a expectativa: deslocar um
\texttt{uint8\_t} por 8 seria \gls{ub} se a operação fosse de 8 bits, e é segura
precisamente por causa da promoção inteira --- a mesma regra que causa o
problema da subseção anterior. A terceira é a armadilha mais frequente do
mundo real: o literal \texttt{1} é \texttt{int}, e o tipo do \emph{destino} não
influencia o tipo da \emph{expressão}. Máscaras de 64 bits exigem
\texttt{1ull}.

\begin{armadilha}[o UBSan pega, o compilador não]
Nenhum dos casos acima gera \emph{warning} com \texttt{-Wall}. O
UndefinedBehaviorSanitizer (\cref{sec:sanitizers}) detecta em tempo de execução
e imprime \texttt{runtime error: shift exponent 32 is too large}. É mais um
argumento para rodar os testes sob \emph{sanitizer} ao menos uma vez.
\end{armadilha}

Vale registrar uma regra que \textbf{mudou}: \texttt{-1 << 1} era \gls{ub} até o
C++17 e passou a ser definido no C++20, porque esse padrão fixou a representação
em complemento de dois como obrigatória.

\subsubsection{\texttt{>>} em \emph{signed} contra \emph{unsigned}}

Em tipos sem sinal, \texttt{>>} insere zeros à esquerda (deslocamento
\emph{lógico}). Em tipos com sinal e valor negativo, insere cópias do bit de
sinal (deslocamento \emph{aritmético}), preservando o valor:

\begin{lstlisting}[style=bt, language=bash,
    caption={Saída de \texttt{src/bits/}: os mesmos bits, dois resultados.}]
int32_t(-16) >> 2 = -4           ← preencheu com o bit de SINAL
uint32_t     >> 2 = 1073741820   ← preencheu com zeros

-7 >> 1 = -4    (arredonda para -∞)
-7 /  2 = -3    (trunca para zero)
\end{lstlisting}

As duas últimas linhas desmentem a substituição que quase todo programador já
fez: \texttt{x >> 1} \textbf{não} é \texttt{x / 2} para negativos.

\begin{regra}[\emph{unsigned} para bits, \emph{signed} para aritmética]
Para manipulação de bits, use sempre tipos sem sinal: eles não têm bit de sinal
para propagar, não têm estouro com comportamento indefinido e não mudam de
significado ao deslocar. Tipos com sinal são para aritmética, onde o sinal
significa alguma coisa.
\end{regra}

\subsubsection{Extrair e inserir campos, de forma portável}

O padrão que aparece em todo protocolo binário e em todo registrador de hardware
é um inteiro que empacota vários campos. Fazê-lo à mão, com deslocamento e
máscara, dá um tipo pequeno e \gls{constexpr}:

\begin{lstlisting}[style=bt, language=C++,
    caption={Transcrição de \texttt{src/bits/demo\_operadores.hpp}: campos com
             layout definido por você.}]
struct CampoBits {
    unsigned deslocamento;
    unsigned largura;

    constexpr std::uint32_t mascara() const {
        return ((1u << largura) - 1u) << deslocamento;
    }
    constexpr std::uint32_t extrair(std::uint32_t reg) const {
        return (reg & mascara()) >> deslocamento;
    }
    constexpr std::uint32_t inserir(std::uint32_t reg, std::uint32_t valor) const {
        return (reg & ~mascara()) | ((valor << deslocamento) & mascara());
    }
};

constexpr CampoBits VERSAO {24, 8};
constexpr CampoBits TIPO   {16, 8};
constexpr CampoBits TAMANHO{ 4, 12};
constexpr CampoBits FLAGS  { 0, 4};
\end{lstlisting}

Como tudo é \texttt{constexpr}, o layout é verificável em compilação, e um
deslocamento errado vira erro de \emph{build}:

\begin{lstlisting}[style=bt, language=C++,
    caption={Transcrição de \texttt{src/bits/demo\_operadores.hpp}: o layout
             conferido pelo compilador.}]
    static_assert(VERSAO.mascara()  == 0xFF000000u);
    static_assert(TAMANHO.mascara() == 0x0000FFF0u);
    static_assert(TAMANHO.extrair(TAMANHO.inserir(0, 4095)) == 4095);
\end{lstlisting}

Guarde este exemplo: ele é a alternativa \textbf{portável} aos campos de bits
nativos, e a diferença entre os dois é o assunto de duas subseções adiante.

\subsubsection{\cabecalho{bit}: as operações que o hardware já sabia fazer}

Todo processador moderno tem instruções para contar bits, achar o bit mais
significativo e rotacionar. Até o C++20 não havia forma portável de chamá-las:
o código usava \texttt{\_\_builtin\_popcount} no GCC,
\texttt{\_BitScanReverse} no MSVC, ou um laço que o compilador talvez
reconhecesse. O \cabecalho{bit} padronizou tudo.

\paragraph{\texttt{bit\_cast}: o fim do \emph{type punning} informal.}
\emph{Type punning} é reinterpretar os bits de um objeto como se fossem de outro
tipo --- operação legítima e comum em serialização e em algoritmos de ponto
flutuante. Todas as formas antigas eram inválidas ou incômodas.

\begin{table}[htbp]
\centering
\caption{As quatro formas de reinterpretar bits, e o veredito de cada uma.}
\label{tab:type-punning}
\footnotesize
\begin{tabular}{@{}ll@{}}
\toprule
\textbf{Forma} & \textbf{Veredito} \\
\midrule
\texttt{*(int*)\&f}                & viola \emph{strict aliasing} --- \gls{ub} \\
\texttt{union \{ float f; int i; \}} & ler o membro inativo é \gls{ub} em C++ \\
\texttt{memcpy(\&i, \&f, 4)}       & correto, mas verboso e \textbf{não} é \gls{constexpr} \\
\texttt{std::bit\_cast<int>(f)}    & correto, conciso \textbf{e} \gls{constexpr} \\
\bottomrule
\end{tabular}
\end{table}

\emph{Strict aliasing} é a regra que permite ao compilador supor que ponteiros
de tipos diferentes não apontam para o mesmo objeto --- suposição sem a qual
quase nenhuma otimização de laço seria possível. Violá-la não produz erro:
produz código que funciona em \texttt{-O0} e falha em \texttt{-O2}.

\texttt{bit\_cast} exige que os dois tipos tenham o mesmo \texttt{sizeof} e
sejam trivialmente copiáveis, e verifica isso em compilação. E, por ser
\gls{constexpr}, a representação de um \texttt{float} passa a ser conferível por
\texttt{static\_assert}:

\begin{lstlisting}[style=bt, language=C++,
    caption={Transcrição de \texttt{src/bits/demo\_bit.hpp}: a representação
             conferida em compilação.}]
    constexpr auto bits_const = std::bit_cast<std::uint32_t>(2.0f);
    static_assert(bits_const == 0x40000000u);
\end{lstlisting}

\paragraph{Contagem.} Três funções, e um uso prático para cada uma:

\begin{lstlisting}[style=bt, language=bash,
    caption={Saída de \texttt{src/bits/}: contagem sobre um valor de 32 bits.}]
v = 00001011011000000000000000000000
  popcount    = 5    bits 1 no total
  countl_zero = 4    zeros à esquerda
  countr_zero = 21   zeros à direita
\end{lstlisting}

\texttt{popcount} dá a distância de Hamming entre dois valores (basta contar os
bits de \texttt{a \textasciicircum{} b}) e a cardinalidade de um conjunto
representado como máscara. \texttt{countl\_zero} dá a posição do bit mais
significativo, e portanto o logaritmo de base 2 inteiro.
\texttt{countr\_zero} dá o alinhamento de um endereço e a posição do próximo
\emph{slot} livre num alocador de bitmap.

\paragraph{Potências de dois.} Quatro funções respondem a perguntas distintas:
\texttt{has\_single\_bit} (``é potência de 2?''), \texttt{bit\_width} (quantos
bits são necessários), \texttt{bit\_floor} (maior potência de 2 menor ou igual)
e \texttt{bit\_ceil} (menor potência de 2 maior ou igual).

\begin{table}[htbp]
\centering
\caption{Saída de \texttt{src/bits/}: as quatro funções sobre dois valores.}
\label{tab:potencias-dois}
\footnotesize
\begin{tabular}{@{}lcccc@{}}
\toprule
\textbf{Valor} & \texttt{has\_single\_bit} & \texttt{bit\_width} & \texttt{bit\_floor} & \texttt{bit\_ceil} \\
\midrule
100  & não & 7  & 64  & 128 \\
1000 & não & 10 & 512 & 1024 \\
\bottomrule
\end{tabular}
\end{table}

O uso canônico é dimensionar um \emph{buffer}: pedir mil elementos e alocar
\texttt{bit\_ceil(1000)}~=~1024 permite trocar o resto de divisão
\texttt{i \% n} pelo mascaramento \texttt{i \& (n-1)}, que é uma instrução em
vez de dezenas. Antes do C++20, \texttt{has\_single\_bit(v)} se escrevia
\texttt{v \&\& !(v \& (v - 1))} --- correto, ilegível, e errado para
\texttt{v == 0} se alguém esquecesse o primeiro termo.

\begin{armadilha}[\texttt{bit\_ceil} tem um caso de \gls{ub}]
\texttt{bit\_ceil(v)} é \gls{ub} se o resultado não couber no tipo. Para
\texttt{uint32\_t}, qualquer \texttt{v} acima de $2^{31}$ pede $2^{32}$, que não
existe no tipo. Verifique antes, ou trabalhe num tipo mais largo.
\end{armadilha}

\paragraph{Rotação.} Rotacionar é diferente de deslocar: os bits que saem de um
lado \textbf{entram} pelo outro. É a operação-base de funções de \emph{hash} e
de criptografia.

\begin{lstlisting}[style=bt, language=bash,
    caption={Saída de \texttt{src/bits/}: deslocamento contra rotação.}]
v          = 10010011
v << 3     = 10011000    (shift: entram zeros, perde bits)
rotl(v, 3) = 10011100    (rotação: os bits VOLTAM)
rotl(v, 11)= 10011100    (11 % 8 = 3 — sem UB)
rotl(v, -3)= 01110010    (negativo inverte o sentido)
\end{lstlisting}

Duas propriedades merecem destaque, porque contrastam com o deslocamento:
\texttt{rotl} reduz a contagem módulo a largura do tipo, e aceita contagem
\textbf{negativa} invertendo o sentido. Não há \gls{ub} para valor nenhum ---
enquanto o idioma manual pré-C++20, \texttt{(v << n) | (v >> (32 - n))}, é
\gls{ub} exatamente quando \texttt{n == 0}, que é o caso que ninguém testa.

\begin{armadilha}[a instrução única depende da \emph{flag} de compilação]
\texttt{popcount} mapeia para a instrução \texttt{POPCNT}, \texttt{countl\_zero}
para \texttt{LZCNT}, \texttt{rotl} para \texttt{ROL} --- mas só se o compilador
tiver permissão para usar a extensão correspondente. Com \texttt{g++ -O2} e
nada mais, \texttt{POPCNT} não é garantido, e o compilador emite a sequência
genérica de cerca de doze instruções. Com \texttt{-msse4.2} ou
\texttt{-march=native}, vira uma. Este subprojeto compila \emph{sem} a
\emph{flag} de propósito: o objetivo é o comportamento, não a instrução.
\end{armadilha}

\subsubsection{\classe{std::bitset}: bits com nome}

\texttt{bitset<N>} é um vetor de $N$ bits empacotados com uma interface de alto
nível: \texttt{b.test(3)} em vez de \texttt{(flags >> 3) \& 1},
\texttt{b.set(3)} em vez de \texttt{flags |= (1u << 3)}, \texttt{b.count()} em
vez de \texttt{\_\_builtin\_popcount(flags)}. É mais legível e igualmente
rápido --- as operações compilam para as mesmas instruções.

Ele suporta as operações de conjunto diretamente: \texttt{|} é união,
\texttt{\&} é interseção, \texttt{\textasciicircum} é diferença simétrica,
\texttt{a \& \textasciitilde b} é diferença, e \texttt{(a \& b).none()} responde
se são disjuntos.

\begin{armadilha}[o índice 0 fica à direita da \emph{string}]
O índice $i$ corresponde ao valor $2^i$, portanto o índice 0 é o bit
\textbf{menos} significativo --- o mais à \textbf{direita} quando se imprime:
\begin{lstlisting}[style=bt, language=bash, frame=none,
                   xleftmargin=0pt, xrightmargin=0pt]
b.set(0) → to_string() = "00000001"
b.set(7) → to_string() = "10000000"
\end{lstlisting}
A regra completa é \texttt{string[j]} $\rightarrow$ índice \texttt{(N-1-j)}. A
\emph{string} lê da esquerda para a direita como um número binário; os índices
contam da direita para a esquerda, como potências. Confundir os dois ao
converter de e para texto é a fonte de \emph{bug} mais comum com
\texttt{bitset}.
\end{armadilha}

Como no \texttt{vector}, há duas formas de acessar com contratos diferentes:
\texttt{test(i)} verifica os limites e lança \texttt{std::out\_of\_range};
\texttt{b[i]} não verifica, e é \gls{ub} fora da faixa. E \texttt{to\_ulong()}
lança \texttt{overflow\_error} se houver bits significativos além do que cabe
num \texttt{unsigned long}.

\begin{table}[htbp]
\centering
\caption{\texttt{bitset} contra as alternativas.}
\label{tab:bitset-alternativas}
\footnotesize
\begin{tabular}{@{}lcccc@{}}
\toprule
\textbf{Opção} & \textbf{Tamanho} & \textbf{Legível} & \textbf{Seguro} & \textbf{Dinâmico} \\
\midrule
\texttt{uint32\_t} + máscara & fixo & não & não & não \\
\texttt{enum class} de \emph{flags} & fixo & \textbf{sim} & \textbf{sim} & não \\
\texttt{std::bitset<N>}      & fixo & \textbf{sim} & sim & não \\
\texttt{std::vector<bool>}   & \textbf{dinâmico} & sim & não$^*$ & \textbf{sim} \\
\texttt{std::vector<char>}   & dinâmico & sim & sim & sim \\
\bottomrule
\end{tabular}\\[2pt]
\footnotesize $^*$ \texttt{operator[]} devolve um \emph{proxy}, não um
\texttt{bool\&} --- ver \cref{sec:stl}.
\end{table}

A comparação com \classe{std::vector} de \texttt{bool} merece ser explícita,
porque os dois empacotam bits e a escolha entre eles não é de gosto.
\texttt{vector<bool>} é a única opção quando o tamanho só é conhecido em tempo
de execução, e paga por isso com o \emph{proxy}: \texttt{auto x = vb[0]} não
produz um \texttt{bool}, produz uma referência disfarçada, e atribuir a ela
modifica o vetor. \texttt{bitset<N>} tem tamanho fixo em compilação e devolve
\texttt{bool} de verdade em \texttt{test()}. Quando o tamanho é dinâmico e o
\emph{proxy} incomoda, a resposta é \texttt{vector<char>} ou
\texttt{boost::dynamic\_bitset}.

O tamanho do \texttt{bitset} é arredondado para múltiplos da palavra:

\begin{lstlisting}[style=bt, language=bash,
    caption={Saída de \texttt{src/bits/}: o arredondamento para a palavra.}]
sizeof: bitset<8>=8  bitset<64>=8  bitset<65>=16  bitset<100>=16
\end{lstlisting}

Um único bit a mais que 64 dobra o custo. Para conjuntos grandes e homogêneos
--- o crivo de Eratóstenes do subprojeto usa \texttt{bitset<1001>}, cerca de 126
bytes --- o ganho é de oito vezes sobre um vetor de \texttt{bool} nativo, e o
que importa não é a memória em si: é que a estrutura inteira cabe no cache L1.

\subsubsection{Campos de bits nativos, e por que evitá-los}

O C++ tem sintaxe própria para empacotar campos num inteiro:
\texttt{unsigned versao : 4;}. Parece a solução ideal para protocolos binários.
\textbf{Não é} --- porque quase tudo sobre o layout é \emph{definido pela
implementação}: a ordem dos campos dentro da palavra, se um campo pode
atravessar a fronteira da unidade de alocação, o alinhamento e o tamanho total.

A ``unidade de alocação'' aqui é a palavra de memória que o compilador escolhe
para acomodar os campos --- e a escolha é dele. O subprojeto \emph{mede} o
layout no GCC para x86-64 e encontra o primeiro campo nos bits menos
significativos; nada garante que outro compilador faça igual.

\begin{armadilha}[um \emph{struct} com campos de bits não é um formato]
Escrever um cabeçalho de protocolo com campos de bits e enviá-lo pela rede é um
\emph{bug} esperando o dia em que a outra ponta usa outro compilador ou outra
arquitetura. O formato precisa ser definido pelo \textbf{seu} código --- é
exatamente o que o \texttt{CampoBits} da subseção anterior faz, ao custo de mais
verbosidade.
\end{armadilha}

\begin{armadilha}[um campo com sinal de 1 bit vale $-1$, não $1$]
A faixa de \texttt{int x : 1} é $\{-1, 0\}$, não $\{0, 1\}$: o único bit
\emph{é} o bit de sinal. Atribuir 1 e ler $-1$ é o comportamento correto, e o
GCC avisa (\texttt{overflow in conversion... changes value from '1' to '-1'}).
O \emph{bug} clássico:
\begin{lstlisting}[style=bt, language=C++, frame=none,
                   xleftmargin=0pt, xrightmargin=0pt]
struct Flags { int ativo : 1; };
f.ativo = 1;
if (f.ativo == 1) ...    // FALSO! f.ativo é -1
if (f.ativo)      ...    // ✓ verdadeiro (-1 é truthy)
\end{lstlisting}
Campos de bits são \textbf{sempre} \texttt{unsigned} ou \texttt{bool}. E
\texttt{char x : 1} tem o mesmo problema, porque a existência de sinal em
\texttt{char} também é definida pela implementação.
\end{armadilha}

\begin{armadilha}[um campo de bits não tem endereço]
Um campo de bits não ocupa bytes inteiros, logo não tem endereço próprio:
\texttt{\&x.campo} não compila, e nenhuma referência pode ser ligada a ele. Isso
significa que passá-lo a qualquer função que receba por referência ---
\texttt{fmt::print}, por exemplo --- falha com \texttt{error: cannot bind
bit-field to 'unsigned char\&'}. O contorno é o \texttt{+} unário, que promove o
campo a um \texttt{int} temporário: \texttt{fmt::print("\{\}", +h.versao)}. O
demo precisou do \texttt{+} em cada campo, e isso \emph{é} a lição.
\end{armadilha}

Há ainda dois custos que não são óbvios. Ler ou escrever um campo vira
deslocamento mais máscara --- mais instruções que acessar um byte inteiro. E
dois campos na mesma palavra \textbf{não} podem ser acessados concorrentemente:
escrever num deles faz leitura-modificação-escrita da palavra inteira, o que é
uma \gls{datarace} pela definição que o C++11 introduziu (\cref{sec:concurrency}).

\begin{table}[htbp]
\centering
\caption{A ferramenta certa para cada objetivo com bits.}
\label{tab:bits-ferramenta}
\footnotesize
\begin{tabular}{@{}ll@{}}
\toprule
\textbf{Objetivo} & \textbf{Ferramenta} \\
\midrule
economizar memória \textbf{interna} ao programa & campos de bits \\
protocolo binário ou arquivo                    & máscaras e deslocamento \\
registrador de hardware                         & máscaras e \texttt{volatile} \\
conjunto de \emph{flags} nomeadas               & \texttt{enum class} com operadores \\
muitos bits homogêneos                          & \classe{std::bitset} \\
\bottomrule
\end{tabular}
\end{table}

\subsubsection{\emph{Endianness}, \texttt{bit\_cast} e serialização}

Um \texttt{uint32\_t} com o valor \texttt{0x12345678} ocupa quatro bytes. A
\textbf{ordem} desses bytes na memória depende da arquitetura: em
\emph{little-endian} (x86, ARM na configuração usual), o byte menos
significativo fica no menor endereço; em \emph{big-endian} (PowerPC, e a
convenção de rede), o mais significativo.

\begin{lstlisting}[style=bt, language=bash,
    caption={Saída de \texttt{src/bits/}: os bytes na memória, nesta máquina.}]
uint32_t 0x12345678 → 78 56 34 12      ← little-endian: menos significativo 1º
\end{lstlisting}

Enquanto os dados vivem na memória do processo, isso é invisível. Assim que eles
viram bytes --- arquivo, \emph{socket}, memória compartilhada --- a ordem passa
a fazer parte do formato. Duas observações fecham o conceito: o \emph{endianness}
afeta a ordem dos \textbf{bytes}, nunca a dos bits dentro de cada byte; e, desde
o C++20, \texttt{std::endian::native} responde a pergunta em tempo de
\emph{compilação}, permitindo \texttt{if constexpr} --- antes disso o truque era
ler o primeiro byte de um \texttt{uint16\_t} por um ponteiro para
\texttt{char}, o que tecnicamente violava \emph{strict aliasing}.

\begin{armadilha}[serializar o \emph{struct} inteiro está sempre errado]
\begin{lstlisting}[style=bt, language=bash, frame=none,
                   xleftmargin=0pt, xrightmargin=0pt]
struct Pacote { uint16_t tipo; uint32_t tamanho; uint8_t versao; };
  sizeof = 12 bytes (os campos somam 7 — o resto é PADDING)

write(fd, &p, sizeof(p)) escreve:
  02 01 00 00 06 05 04 03 07 00 00 00
\end{lstlisting}
São três problemas, todos silenciosos. \textbf{Padding}: os bytes que o
compilador insere entre campos para respeitar o alinhamento têm valor
indefinido, e enviá-los vaza memória do processo --- já foi CVE.
\textbf{\emph{Endianness}}: a ordem gravada é a da arquitetura de quem escreveu,
não a do formato. \textbf{Layout}: outro compilador pode ordenar ou alinhar
diferente. Nenhum dos três aparece em teste, porque quem escreve e quem lê é o
\emph{mesmo binário}.
\end{armadilha}

A forma correta é serializar campo a campo, numa ordem que o \emph{seu} código
define:

\begin{lstlisting}[style=bt, language=C++,
    caption={Transcrição de \texttt{src/bits/demo\_endianness.hpp}: escrita em
             \emph{big-endian}, independente da máquina.}]
    template <typename T>
    void escrever(T v) {
        static_assert(std::is_unsigned_v<T>, "serialize só tipos sem sinal");
        for (int i = sizeof(T) - 1; i >= 0; --i) {
            saida_.push_back(static_cast<std::uint8_t>(v >> (i * 8)));
        }
    }
\end{lstlisting}

\begin{lstlisting}[style=bt, language=bash,
    caption={Saída de \texttt{src/bits/}: sete bytes contra doze, e ida e volta
             exata.}]
serializado campo a campo, big-endian:
  01 02 03 04 05 06 07        ← 7 bytes, contra 12 do struct cru
round-trip: ✓ idêntico
\end{lstlisting}

\begin{porque}[por que o deslocamento resolve o \emph{endianness}]
Porque \texttt{v >> (i*8)} é uma operação sobre o \textbf{valor numérico}, não
sobre os bytes. \texttt{v >> 24} devolve o byte mais significativo em qualquer
arquitetura, porque ``mais significativo'' é uma propriedade do número, não da
memória. \emph{Endianness} só existe quando alguém olha os bytes --- e este
código nunca olha. É o mecanismo por trás de Protobuf, MessagePack e Cap'n
Proto.
\end{porque}

Quando é mesmo necessário inverter a ordem de um valor inteiro --- ao ler um
formato de outra convenção, por exemplo --- a ferramenta é
\texttt{std::byteswap}, do C++23, que vira uma instrução \texttt{BSWAP}. As
alternativas, em ordem decrescente de qualidade, são
\texttt{\_\_builtin\_bswapNN} do GCC e do Clang, o par
\texttt{htonl}/\texttt{ntohl} do POSIX (limitado a 16 e 32 bits) e os
deslocamentos escritos à mão, que o compilador em geral reconhece e converte na
mesma instrução.

% =============================================================================
\subsection{O IO de arquivo com \emph{streams}}
\label{sec:io-streams}
% =============================================================================

A \cref{sec:file-io} apresentou a hierarquia dos \emph{streams} e o
\classe{std::filesystem}; esta seção completa o assunto pelo lado prático, com o
que o subprojeto \projeto{file\_io} demonstra em código.

Um \emph{stream} é uma abstração de sequência de caracteres com direção. Por
baixo dele há um \texttt{streambuf}, o objeto que efetivamente conversa com o
destino --- um arquivo, um terminal, uma região de memória --- e que guarda o
\emph{buffer}: uma área intermediária onde os caracteres se acumulam antes de
serem entregues, para que uma única chamada ao sistema operacional escreva
muitos deles de uma vez.

\subsubsection{Os quatro objetos globais, e a diferença que importa}

\texttt{std::cin} lê da entrada padrão; \texttt{std::cout} escreve na saída
padrão, com \emph{buffer}; \texttt{std::cerr} e \texttt{std::clog} escrevem
ambos no erro padrão, e a diferença entre eles é exatamente o \emph{buffer}.

\begin{regra}[\texttt{cerr} para o que precisa sobreviver ao \emph{crash}]
\texttt{std::cerr} não tem \emph{buffer}: cada escrita chega ao destino
imediatamente. Isso custa caro --- uma chamada ao sistema por operação --- e é
precisamente o que se quer da última mensagem antes de um aborto.
\texttt{std::clog} escreve no mesmo lugar com \emph{buffer}, e é a escolha certa
para \emph{log} volumoso, desde que se aceite perder as últimas linhas num
encerramento anormal. Pela mesma razão, prefira \texttt{'\textbackslash n'} a
\texttt{std::endl}: o segundo é o primeiro \emph{mais um flush}, e o flush
desnecessário em laço é um dos custos de IO mais comuns.
\end{regra}

\subsubsection{Manipuladores: poderosos e pegajosos}

Os manipuladores de \cabecalho{iomanip} controlam base (\texttt{dec},
\texttt{hex}, \texttt{oct}), notação de ponto flutuante (\texttt{fixed},
\texttt{scientific}), precisão, largura, preenchimento e alinhamento.

\begin{armadilha}[o manipulador muda o \emph{stream}, não a expressão]
\texttt{std::hex} não formata o próximo valor: ele altera o \textbf{estado} do
\emph{stream}, e o estado permanece. Um \texttt{std::hex} esquecido faz toda a
saída seguinte sair em hexadecimal, inclusive a de outra função, de outro
arquivo, escrita por outra pessoa. A exceção é \texttt{std::setw}, que vale
\emph{apenas} para a próxima inserção --- o que torna o conjunto ainda menos
previsível. É a razão de existir \texttt{std::format} (\cref{sec:format}), cuja
especificação é local à chamada e não deixa resíduo.
\end{armadilha}

\subsubsection{Arquivos: texto, binário e posicionamento}

As três classes de \cabecalho{fstream} correspondem às três direções:
\texttt{ofstream} escreve, \texttt{ifstream} lê, \texttt{fstream} faz as duas.
Todas são \gls{raii}: o destrutor fecha o arquivo e esvazia o \emph{buffer}.

\begin{table}[htbp]
\centering
\caption{Os modos de abertura mais usados.}
\label{tab:modos-abertura}
\footnotesize
\begin{tabular}{@{}ll@{}}
\toprule
\textbf{Declaração} & \textbf{Efeito} \\
\midrule
\texttt{ofstream f("a.txt");}                 & \texttt{out | trunc} --- apaga o conteúdo \\
\texttt{ofstream f("a.txt", ios::app);}       & acrescenta ao fim \\
\texttt{ifstream f("a.txt");}                 & \texttt{in} \\
\texttt{fstream f("a.txt", ios::in|ios::out);}& lê e escreve, sem truncar \\
\texttt{ofstream f("a.bin", ios::binary);}    & binário, sem tradução de fim de linha \\
\bottomrule
\end{tabular}
\end{table}

O posicionamento tem \textbf{dois} cursores independentes: \texttt{tellg} e
\texttt{seekg} governam a leitura (\emph{get}), \texttt{tellp} e \texttt{seekp}
governam a escrita (\emph{put}). As origens são \texttt{ios::beg},
\texttt{ios::cur} e \texttt{ios::end}. É com eles que se lê um arquivo inteiro
de uma vez --- ir ao fim, medir a posição, voltar ao início, ler tudo --- e que
se faz acesso aleatório a registros de tamanho fixo, saltando
\texttt{n * sizeof(Registro)} bytes.

\begin{armadilha}[o modo binário do demo é um exemplo do que não fazer em rede]
O subprojeto grava \emph{structs} com \texttt{ofs.write(reinterpret\_cast<const
char*>(\&r), sizeof(Registro))}, e isso é legítimo enquanto o arquivo é privado
do programa. Para qualquer formato que atravesse máquinas, compiladores ou
versões, vale integralmente a advertência da subseção anterior: \emph{padding},
\emph{endianness} e layout tornam o resultado dependente de quem escreveu.
Formato que sai do processo se serializa campo a campo.
\end{armadilha}

Quanto a erros, um \emph{stream} tem quatro estados consultáveis ---
\texttt{good}, \texttt{fail}, \texttt{bad}, \texttt{eof} --- e, por padrão,
\textbf{não lança}: uma leitura que falha apenas liga um bit, e é o
\texttt{if (!ifs)} que precisa existir. Quem preferir exceções pode pedi-las com
\texttt{ifs.exceptions(std::ios::failbit | std::ios::badbit)}, e aí um
\texttt{open} malsucedido lança \texttt{std::ios\_base::failure}.

\subsubsection{\emph{String streams}: a mesma interface, em memória}

\texttt{ostringstream} monta texto, \texttt{istringstream} \emph{parseia},
\texttt{stringstream} faz as duas sobre o mesmo \emph{buffer}. O uso mais
frequente do segundo é a tokenização: \texttt{std::getline(iss, token, ',')}
separa um CSV, e \texttt{std::getline(iss, chave, '=')} seguido de
\texttt{std::getline(iss, valor)} \emph{parseia} uma linha de configuração.

Duas notas de desempenho. Primeiro, \texttt{str()} devolve uma \textbf{cópia} da
\emph{string} interna; desde o C++20 existe \texttt{view()}, que devolve um
\classe{std::string\_view} sem alocar --- e que, como todo \emph{view}, é
invalidado assim que o \emph{stream} é modificado. Segundo, para converter texto
em número, \emph{string streams} são a opção mais cara das cinco da
\cref{tab:texto-numero}: cada conversão constrói um objeto de \emph{stream}
inteiro, com \emph{locale} e \emph{flags}. Para dados, \cabecalho{charconv}.

\subsubsection{\texttt{rdbuf}: o ponto de extensão que torna tudo componível}

O método \texttt{rdbuf()} devolve --- e, com argumento, \textbf{troca} --- o
\texttt{streambuf} de um \emph{stream}. É isso que permite redirecionar
\texttt{std::cout} para um arquivo sem tocar em nenhuma das chamadas que
escrevem nele, e capturar a saída de uma função em uma \emph{string} para
testá-la.

Restaurar o \emph{buffer} à mão é frágil --- uma exceção no meio deixa
\texttt{cout} apontando para um arquivo já destruído --- e é exatamente o
problema que \gls{raii} resolve:

\begin{lstlisting}[style=bt, language=C++,
    caption={Transcrição de \texttt{src/file\_io/demo\_logging.hpp}: o
             redirecionamento que se desfaz sozinho.}]
class RedirectGuard {
public:
    // Redireciona 'stream' para 'destino'
    explicit RedirectGuard(std::ostream& stream, std::ostream& destino)
        : stream_(stream), buf_original_(stream.rdbuf(destino.rdbuf())) {}

    ~RedirectGuard() {
        stream_.rdbuf(buf_original_);   // restaura automaticamente
    }

    RedirectGuard(const RedirectGuard&)            = delete;
    RedirectGuard& operator=(const RedirectGuard&) = delete;

private:
    std::ostream&   stream_;
    std::ofstream   arquivo_;        // vazio se redirecionou para outro stream
    std::streambuf* buf_original_;
};
\end{lstlisting}

E, uma vez que o \texttt{streambuf} é um ponto de extensão, escrever um novo dá
capacidades que nenhuma \emph{flag} oferece. O \texttt{tee} do Unix, que envia a
mesma saída ao terminal e a um arquivo, cabe em duas funções virtuais:

\begin{lstlisting}[style=bt, language=C++,
    caption={Transcrição de \texttt{src/file\_io/demo\_logging.hpp}: um
             \texttt{streambuf} que escreve em dois destinos.}]
class TeeBuf : public std::streambuf {
public:
    TeeBuf(std::streambuf* a, std::streambuf* b) : a_(a), b_(b) {}

protected:
    int overflow(int c) override {
        if (c == EOF) return !EOF;
        if (a_->sputc(static_cast<char>(c)) == EOF) return EOF;
        if (b_->sputc(static_cast<char>(c)) == EOF) return EOF;
        return c;
    }

    std::streamsize xsputn(const char* s, std::streamsize n) override {
        a_->sputn(s, n);
        b_->sputn(s, n);
        return n;
    }

private:
    std::streambuf* a_;
    std::streambuf* b_;
};
\end{lstlisting}

\texttt{overflow} é chamado quando o \emph{buffer} enche ou quando um caractere
solto precisa ser entregue; \texttt{xsputn} é o caminho rápido para blocos.
Sobrescrever os dois é o contrato mínimo de um \texttt{streambuf} de saída.

\subsubsection{Um \emph{logger}, e o que ele reúne}

O \texttt{Logger} do subprojeto junta quase tudo o que o capítulo tratou:
carimbo de tempo do \cabecalho{chrono}, formatação, escrita simultânea em
terminal e arquivo, e um \texttt{std::mutex} para que várias \emph{threads}
possam registrar sem entrelaçar linhas (\cref{sec:concurrency}).

\begin{lstlisting}[style=bt, language=bash,
    caption={Saída de \texttt{src/file\_io/}: níveis, carimbo e destino.}]
[08:42:01.123] [TRACE] inicializando...
[08:42:01.124] [DEBUG] variável x = 42
[08:42:01.124] [INFO ] conexão estabelecida
[08:42:01.125] [WARN ] buffer 80% cheio
[08:42:01.125] [ERROR] falha ao abrir arquivo   ← vai para stderr
[08:42:01.125] [FATAL] estado inválido          ← vai para stderr
\end{lstlisting}

Duas decisões de projeto merecem nota. A primeira é o roteamento por nível:
\texttt{ERROR} e \texttt{FATAL} vão para \texttt{stderr}, o resto para
\texttt{stdout}, de modo que \texttt{./prog > saida.txt} deixa os erros
visíveis no terminal. A segunda é o \texttt{flush} explícito após cada linha
gravada no arquivo: ele custa caro e é a única forma de garantir que a última
mensagem antes de um \emph{crash} chegue ao disco --- a mesma troca de
\texttt{cerr} contra \texttt{clog}, agora sob controle do programa.

\begin{table}[htbp]
\centering
\caption{Redirecionamento no \emph{shell} --- a referência que acompanha o
         subprojeto.}
\label{tab:redirecionamento}
\footnotesize
\begin{tabular}{@{}ll@{}}
\toprule
\textbf{Comando} & \textbf{Efeito} \\
\midrule
\texttt{./prog > out.txt}            & \emph{stdout} para arquivo (trunca) \\
\texttt{./prog >> out.txt}           & \emph{stdout} para arquivo (acrescenta) \\
\texttt{./prog 2> err.txt}           & \emph{stderr} para arquivo \\
\texttt{./prog > out.txt 2>\&1}      & \emph{stderr} segue \emph{stdout} \\
\texttt{./prog 2>\&1 | tee log.txt}  & terminal e arquivo ao mesmo tempo \\
\texttt{./prog > /dev/null 2>\&1}    & descarta tudo \\
\bottomrule
\end{tabular}
\end{table}

\begin{armadilha}[a ordem do redirecionamento importa]
\texttt{./prog > out.txt 2>\&1} e \texttt{./prog 2>\&1 > out.txt} não são a
mesma coisa. O \texttt{2>\&1} copia o destino \emph{atual} do descritor 1; no
segundo comando ele é executado \textbf{antes} de o descritor 1 ser
redirecionado, então \emph{stderr} continua no terminal e só \emph{stdout} vai
para o arquivo. Os descritores são 0 para \emph{stdin}, 1 para \emph{stdout} e 2
para \emph{stderr}.
\end{armadilha}

\subsubsection{\classe{std::filesystem}, em uma linha por conceito}

Para completar: \texttt{fs::path} compõe caminhos com \texttt{operator/} e os
decompõe com \texttt{stem()}, \texttt{extension()}, \texttt{parent\_path()} e
\texttt{filename()}, além de normalizar com \texttt{lexically\_normal()};
\texttt{create\_directories} cria a hierarquia inteira;
\texttt{directory\_iterator} e \texttt{recursive\_directory\_iterator}
percorrem diretórios como
\emph{ranges} (\cref{sec:ranges}); \texttt{file\_size},
\texttt{last\_write\_time}, \texttt{status().permissions()} e \texttt{space}
dão os metadados.

E, como registrado na \cref{sec:file-io}, cada função tem duas assinaturas: uma
que lança \texttt{fs::filesystem\_error} e outra que recebe um
\texttt{std::error\_code\&} e nunca lança. O critério de escolha é a
\emph{frequência esperada} do erro --- um arquivo de configuração que deveria
existir e não existe é excepcional; um arquivo que some no meio de uma varredura
de diretório não é.

% =============================================================================
\subsection{Síntese}
% =============================================================================

Há um padrão atravessando os quatro subprojetos deste capítulo, e vale nomeá-lo:
\textbf{o custo está sempre onde a abstração é mais confortável}.
\classe{std::string} é confortável porque possui os dados, e é por possuí-los
que ela aloca; \classe{std::string\_view} custa zero porque não possui nada, e é
por não possuir que ela pende. \texttt{std::regex} aceita qualquer padrão, e é
por aceitar qualquer padrão que ela não pode ser rápida para um padrão fixo. Os
campos de bits nativos escondem os deslocamentos, e é por escondê-los que
deixam o layout com o compilador. Os \emph{iostreams} guardam estado de
formatação, e é por guardá-lo que o \texttt{std::hex} esquecido contamina a
saída inteira.

Em todos os casos, o preço da propriedade desejada é a armadilha que vem junto.

O segundo padrão é metodológico, e é o que este capítulo tem de mais
transferível. Em quatro ocasiões a medição contrariou quem a fez, e nas quatro a
lição estava no desacordo, não no número:

\begin{itemize}
\item o \emph{benchmark} de formatação \textbf{não} mostrou a \texttt{fmt}
      ganhando --- porque o repositório linka a variante de depuração da
      biblioteca, e nenhuma \emph{flag} no código do usuário conserta isso;
\item a \texttt{regex} construída dentro do laço levou mais tempo fazendo
      \textbf{dez vezes menos trabalho}, o que separa o custo de construir do
      custo de casar;
\item o \texttt{static\_assert} sobre o tamanho de um literal recusou 31 e
      exigiu 33, aplicando ao autor a lição de que \texttt{size()} conta bytes;
\item \texttt{high\_resolution\_clock}, cujo nome promete o melhor relógio para
      medir, revelou-se o \textbf{não monotônico} nesta implementação.
\end{itemize}

\begin{regra}[o que sobra quando o número some]
Quatro hábitos resumem todo o capítulo. Conte alocações em vez de cronometrar
quando puder --- é determinístico. Cronometre com \texttt{steady\_clock},
descartando a primeira execução, repetindo e tomando a mediana. Declare as
\emph{flags} de compilação ao lado de qualquer número, incluindo as da
biblioteca de terceiros. E, quando o resultado contrariar a expectativa,
suspeite primeiro da medição --- e, se ela resistir, publique a surpresa: é ela
que ensina.
\end{regra}
